% !TeX program = xelatex
\documentclass[11pt,a4paper]{article}
\usepackage[margin=23mm,headheight=16pt]{geometry}
\usepackage{fontspec}
\IfFontExistsTF{Noto Sans CJK SC}{
\setmainfont{Noto Sans CJK SC}[AutoFakeBold=1.5,AutoFakeSlant=.15]
\setsansfont{Noto Sans CJK SC}[AutoFakeBold=1.5]
}{
\setmainfont{NotoSansCJKsc-Regular.otf}[Path=fonts/,AutoFakeBold=1.5,AutoFakeSlant=.15]
\setsansfont{NotoSansCJKsc-Regular.otf}[Path=fonts/,AutoFakeBold=1.5]
}
\XeTeXlinebreaklocale "zh"
\XeTeXlinebreakskip=0pt plus 1pt
\usepackage{amsmath,amssymb,amsthm,bm,mathtools}
\usepackage{booktabs,tabularx,array,graphicx,xcolor}
\usepackage{xurl,hyperref,fancyhdr,enumitem}
\usepackage{tikz}
\usetikzlibrary{arrows.meta,positioning}
\usepackage{needspace}
\hypersetup{colorlinks=true,linkcolor=blue!45!black,urlcolor=blue!45!black,citecolor=blue!45!black,pdftitle={Anatomical GPS V4.3: 材料技能与最小任务校准}}
\setlength{\emergencystretch}{4em}
\setlength{\parskip}{.35em}
\setlength{\parindent}{1.4em}
\linespread{1.06}
\setlist{nosep,leftmargin=2em}
\newcolumntype{Y}{>{\raggedright\arraybackslash}X}
\numberwithin{equation}{section}
\newtheorem{theorem}{定理}[section]
\newtheorem{proposition}[theorem]{命题}
\newtheorem{lemma}[theorem]{引理}
\newtheorem{corollary}[theorem]{推论}
\theoremstyle{definition}
\newtheorem{definition}[theorem]{定义}
\newtheorem{assumption}[theorem]{假设}
\newtheorem{remark}[theorem]{说明}
\renewcommand{\proofname}{证明}
\renewcommand{\abstractname}{摘要}
\renewcommand{\refname}{参考文献}
\renewcommand{\contentsname}{目录}
\DeclareMathOperator{\rank}{rank}
\DeclareMathOperator{\range}{range}
\DeclareMathOperator{\diag}{diag}
\newcommand{\R}{\mathbb R}
\newcommand{\norm}[1]{\left\lVert#1\right\rVert}
\pagestyle{fancy}\fancyhf{}
\fancyhead[L]{\small Anatomical GPS V4.3 · 完整理论研究稿}
\fancyhead[R]{\small 材料技能与最小任务校准}
\fancyfoot[C]{\thepage}
\begin{document}
\title{Anatomical GPS V4.3\\[2mm]\Large 材料几何 GP 技能与最小反馈校准\\[2mm]\large 任务保持、允许实现选择、可执行探查与二阶恢复}
\author{完整公式、算法接口与独立审读修订稿}
\date{2026年10月9日}
\maketitle
\begin{abstract}
本文将迁移的单位定义为材料任务及其允许的机器人实现，而非固定的世界笛卡尔轨迹。对局部结构化响应 $J=H+\Theta N$，将任务零空间与实际校准信息共同商掉，得到不可消除耦合谱，给出零残差历史上任意共同首次物理校正动作的精确最小最大误差、谱预算前沿和精确任务响应的必要充分方向数。实际探查进一步产生凸修复、输入能力与二阶任务恢复证书。技能实现唯一采用 GP 学习 forcing：体表材料切空间中的相对偏置和 $SO(3)$ 姿态残余经过明确运输学习，固定几何二阶系统按材料进度推进。调和圆柱图册、SMPL-X 附着与姿态纤维保持其接口；可微逆可达性在允许集合中提议真实可执行实现。名义技能、实际图像纠偏、观测失锁及接触控制分别定义，校准质量同时约束恢复率和允许推进速度。本文保留 4.2 的核心研究主张并补全技能与执行接口。已完成可复现数学、坐标和刚体几何核验；尚未报告真实机器人或真人扫描效果。GP、AL-DMP 与几何 DMP 是已有基础，核心原创性集中于实现选择与必要反馈信息之间的可计算关系。
\end{abstract}
\noindent\textbf{研究主张：}在材料任务及观测质量条件保持时，选择允许的物理实现以降低必须重新校准的反馈方向；以实际可执行探查补足必要方向，并验收修复后的任务恢复性能。
\newpage
\begingroup
\footnotesize
\setlength{\parskip}{0pt}
\tableofcontents
\endgroup
\newpage

\section{需要保持的是任务，而非一个唯一探头姿态}
令 $\mathcal M$ 为有效材料体积图册，任务为 $\mathcal T=(\tau,\mathcal A_\tau,\mathcal Q_\tau)$：$\tau$ 指定血管、器官、材料片段及扫描目的；$\mathcal A_\tau$ 是允许的探头位置、姿态、相位偏差和机器人实现；$\mathcal Q_\tau$ 指定可信目标关联、视野、声窗与特征质量。允许关系必须在优化前定义。若任务要求图像横向坐标严格为零，则把目标改成离中心点不是合法任务保持；若任务只要求同一片段与合格切面、允许横向偏置，则不同偏置可以是等价实现。

对每个实现 $a\in\mathcal A_\tau$，局部可信特征误差 $e\in\R^m$ 描述完成条件；维数 $m$ 必须保持。选取新的 $a$ 可以改变特征参考值，但不能删掉难控制的特征后声称降低了校准维数。下文的保证首先是局部首次单位响应或统一收缩合同；完整长扫描仍需沿路径反复验收。

\subsection{全稿关键符号及单位约定}
\begin{center}\small
\begin{tabularx}{\linewidth}{lY}
\toprule
符号 & 含义与责任\\
\midrule
$\xi=(\varphi,h,d)$、$P,Q$ & 调和材料坐标、身体位置映射与局部姿态标架\\
$s,\nu$、$b,z_b,\Upsilon,\omega$ & 材料进度与推进速度；名义偏置及相对旋转生成状态\\
$H,N,\Theta,E_{\rm geo}$ & 已知任务通道、不确定通道、患者响应及结构外误差\\
$\delta,\eta,\varepsilon_{\rm data}$ & 响应范数界、几何余项界及实际探查数据噪声界\\
$B,D_\star,D_Q$ & 任务实现、不可消除曝光与数据条件剩余曝光\\
$X,\widehat J,C,q$ & 实际激励、估计响应、反馈修复与验收误差\\
$e,\zeta,\omega,S$ & 真实任务误差、纠偏速度、反馈带宽及增广状态包络\\
\bottomrule
\end{tabularx}\end{center}
生成器的姿态空间角速度 $\omega(s)$ 与反馈带宽 $\omega>0$ 在各自章节局部使用；需要同时运算时分别记为 $\omega_{\rm sk}$ 与 $\omega_{\rm fb}$。探查积分运动与瞬时速度使用相应时间尺度，毫米、弧度及任务容差通过固定物理度量归一化。学习空间中的 GP 输出是虚拟形状驱动，不是患者端口的物理接触力。

\section{材料调和圆柱图册、SMPL-X 与姿态纤维的接口}
\subsection{保留圆柱设计，并区分图册和内部解剖}
在每个有效环壳材料子域 $\Omega_a$ 上，保留
\begin{equation}
 \xi=(\varphi,h,d),\qquad
 \Delta_{g_0}h=0,\quad \Delta_{g_0}d=0,\quad
 \Delta_{g_0}\widetilde\varphi=0.
\end{equation}
轴向与径向场取各自端面或边界上的 Dirichlet 条件；角度场在切开的域上取两岸相差 $2\pi$ 的周期 lift，跨缝使用转移图。圆柱核心退化域需要独立图，腿和躯干可以是不同圆柱图；不强行制造一个覆盖整个人体的单图。调和性不保证三维坐标无折叠，必须验收 $\det D_\xi F_{a,0}\ne0$、方向一致性及自交。

SMPL-X 表面采用标准蒙皮形式\cite{pavlakos2019}
\begin{align}
 \bar v_i(c)&=\bar v_i+B_{S,i}(\beta)+B_{E,i}(\psi)+B_{P,i}(\vartheta),\\
 v_i(c)&=\sum_k w_{ik}\bigl(R_k(c)\bar v_i(c)+t_k(c)\bigr).
\end{align}
这里 $c$ 含体型、姿态、表情和世界位姿，骨架变换含静止姿态校正。体积内部需要另外给出权重 $w_k(p_0)$ 和修正场 $b_c(p_0)$；表面 SMPL-X 不会自动提供它们，也不会给出患者内部血管的准确位置。体积延拓与接触形变为
\begin{align}
 L_c(p_0)&=\sum_k w_k(p_0)\bigl[R_{k,c}(p_0+b_c(p_0))+t_{k,c}\bigr],\\
 P_{a,c,f}(\xi)&=C_{c,f}\bigl(L_c(F_{a,0}(\xi))\bigr).
\end{align}
已知身体运输只保持材料地址。内部滑移、压迫、注册误差与目标形态差异仍进入独立未知量或误差界。

\subsection{调和圆柱的边界条件、缝合与跨身体运输}
\label{sec:v43:harmonic-detail}
圆柱坐标是对材料拓扑的建模选择，而不是把人体强行拟合成几何圆柱。对腿、躯干分别取参考材料环壳 $\Omega_a$，设其轴向边界为 $\Gamma_h^-,\Gamma_h^+$，外表面与内参考壳为 $\Gamma_d^0,\Gamma_d^1$。一种明确的场构造是
\begin{align}
 \Delta_{g_0}h&=0,& h|_{\Gamma_h^-}&=0,&h|_{\Gamma_h^+}&=1,\\
 \Delta_{g_0}d&=0,& d|_{\Gamma_d^0}&=0,&d|_{\Gamma_d^1}&=1.
\end{align}
未指定边界使用与该域一致的无通量条件。角度不能作为一个全局单值实函数绕圆柱递增；在切口两岸指定 $\widetilde\varphi^+-\widetilde\varphi^-=2\pi$，并匹配通量，再将 $\varphi=\widetilde\varphi\bmod 2\pi$ 视为 $S^1$ 值场。有效局部图为 $U_a\subset S^1\times(0,1)\times(0,1)$。这保留轴向、环向及内部层次的材料语义，同时避免跨角度接缝的数值大跳变。

这里所谓拓扑保持，指有效图册及其转移映射不改变所覆盖材料域的对应关系。它不等于在任意组织变形、压迫或人体个体差异下都有准确的内部同胚。坐标映射 $F_{a,0}:U_a\to\Omega_a$ 必须局部可逆；环壳内核、腋下、腹股沟及部位连接处可以需要额外图。腿图和躯干图的重叠处若有真实对应，定义 $\xi_b=\psi_{ba}(\xi_a)$；互不重叠的扫描部位保留部位标签 $a$，不虚构一条图间转换。图间同一几何量的表示满足
\begin{equation}
 F_{b,0}\circ\psi_{ba}=F_{a,0},\qquad
 g_b=(D\psi_{ba})^{-\top}g_a(D\psi_{ba})^{-1}.
\end{equation}
本文使用的材料技能尺度是事先固定的参考度量 $g_0$ 或固定任务容差度量；新身体的物理速度与能力由映射的真实 Jacobian 重建。不能把每个身体中的世界欧氏距离都暗中当作同一材料度量，然后声称严格迁移不变。

\subsection{完整 SMPL-X 蒙皮公式及内部延拓的责任边界}
\label{sec:v43:smpl-detail}
为使体表对应可实现，进一步明确前述简写。以 $N_v$ 个顶点的静止网格 $\bar V$、形状参数 $\beta$、表情参数 $\psi$ 和关节旋转参数 $\vartheta$ 定义
\begin{equation}
 T(\beta,\psi,\vartheta)=\bar V+B_S(\beta)+B_E(\psi)+B_P(\vartheta),
 \quad B_S=\sum_l\beta_l S_l,\quad B_E=\sum_j\psi_j E_j.
\end{equation}
姿态修正可由旋转偏离静止姿态的元素构造，例如
\begin{equation}
 B_P(\vartheta)=\sum_{k,j}P_{kj}\,[\operatorname{vec}(R_k(\vartheta)-I)]_j.
\end{equation}
关节位置由模型提供的关节回归器及形状得到 $J(\beta)$。沿运动学树的齐次变换为
\begin{equation}
 G_k(\vartheta,J)=\prod_{j\in\operatorname{anc}(k)}
 \begin{bmatrix}R_j(\vartheta)&J_j-J_{\operatorname{par}(j)}\\0&1\end{bmatrix},
 \qquad \widehat G_k=G_k(\vartheta,J)G_k(0,J)^{-1}.
\end{equation}
根节点按其世界位姿处理。线性混合蒙皮输出
\begin{equation}
 \begin{bmatrix}v_i\\1\end{bmatrix}
 =\sum_k w_{ik}\widehat G_k
 \begin{bmatrix}T_i(\beta,\psi,\vartheta)\\1\end{bmatrix},
 \qquad \sum_k w_{ik}=1.
\end{equation}
这个表达显示了 $R_k,t_k$ 中必须包含静止骨架校正的原因。跨身体的同顶点、同面片及重心坐标提供体表材料对应；它们本身不观测内部结构，也不规定肌肉、血管和脏器随皮肤一起运动。

体积延拓可以在四面体域中插值权重，也可以解弹性或调和延拓问题，但必须作为独立实现记录。接触映射 $C_{c,f}$ 若只由近似组织模型给出，误差进入 $E_{\rm geo}$ 及其可信界。绝不能从表面 SMPL-X 的拟合误差直接推导内部血管误差的大小。使用实际 CoP、人体运动补偿和独立观测校准，是这一边界下保持理论有效的条件。

\subsection{实际动作的材料拉回必须扣除人体运动}
记 $J_P=D_\xi P_{a,c,f}$，则
\begin{align}
 \dot p&=J_P\dot\xi+b_p,
 &b_p&=D_cP\dot c+D_fP\dot f,\\
 \dot\xi&=J_P^{-1}(v_{\rm CoP}-b_p).
\end{align}
使用实际接触作用点 CoP 的运动，不能把位于体外的 TCP 硬塞进体积逆映射。由此校准数据中的输入来自实际动作和身体运动补偿，而非计划中的命令。

姿态纤维取 $R=Q_{a,c,f}(\xi)\rho$，$\rho\in SO(3)$。定义
\begin{equation}
 A_i=(Q^\top\partial_iQ)^\vee,\qquad
 a_c=(Q^\top(D_cQ\dot c+D_fQ\dot f))^\vee .
\end{equation}
若 $\omega_R=(R^\top\dot R)^\vee$，则纤维的相对角速度为
\begin{equation}
 \omega_\rho=\omega_R-\rho^\top(A_\xi\dot\xi+a_c).
\end{equation}
这使探查的平移、姿态和材料位置能够共同记录。更换局部标架 $Q'=QS$ 时，$\rho'=S^\top\rho$，实际姿态仍为同一个 $R$。

\subsection{校准谱必须在物理度量下计算}
输入、任务、不确定通道分别赋正定度量 $W,M,S$。对下文矩阵定义
\begin{equation}
 \widetilde H=M^{1/2}HW^{-1/2},\quad
 \widetilde N=S^{1/2}NW^{-1/2},\quad
 \widetilde\Theta=M^{1/2}\Theta S^{-1/2}.
\end{equation}
全文伪逆、正交投影和奇异值都在这个白化模型中理解。为简洁，后文省略波浪号。度量需要固定到真实的毫米、弧度、任务容差与通道尺度，不能为了获得收益而重设。

被动换坐标 $u'=Tu$, $e'=Pe$, $z'=Vz$ 给出
\begin{equation}
 H'=PHT^{-1},\quad N'=VNT^{-1},\quad
 \Theta'=P\Theta V^{-1},\quad B'=TBP^{-1}.
\end{equation}
度量同步作协变变换，例如 $W'=T^{-\top}WT^{-1}$。白化前后算子由正交基变换关联，故诱导奇异值与任务下界不变。\textbf{重画圆柱坐标不能省校准；真实改变接触姿态、输入分配或传感支持才可能省校准。}

若接触的线性约束是 $E_cu=0$，先取 $\ker E_c$ 的正交基 $Z_c$，将输入写成 $u=Z_c\bar u$，对 $HZ_c,NZ_c$ 推导。实际能力多面体、关节约束和碰撞仍需后续验收。

\subsection{坐标设计服务的理论对象是一个任务耦合商空间}
在材料位置、探头姿态和实现组成的基点 $a$ 上，取允许命令空间 $\mathcal V_a$、任务反馈空间 $\mathcal Y_a$ 与不确定通道空间 $\mathcal Z_a$，映射为
\begin{equation}
 H_a:\mathcal V_a\longrightarrow\mathcal Y_a,\qquad
 N_a:\mathcal V_a\longrightarrow\mathcal Z_a.
\end{equation}
任务零空间可以抵消的通道是 $\mathcal C_a=N_a\ker H_a$。将它商掉，定义与右逆选择无关的诱导映射
\begin{equation}
 \overline D_a:\mathcal Y_a\longrightarrow\mathcal Z_a/\mathcal C_a,
 \qquad \overline D_av=[N_aB_av],\quad H_aB_a=I.
\end{equation}
任何两个右逆之差都在 $\ker H_a$ 中，因此上述等价类相同。在物理度量下选取正交代表，就得到后文 $D_\star$。完成探查后的信息子空间 $Q_a$ 再被商掉，对应 $\mathcal Z_a/(\mathcal C_a+Q_a)$ 和后文 $D_Q$。

这才是调和图册与纤维的理论用途：在各人体、各图和姿态标架中记录同一个任务与通道等价类；允许不同机器人用不同动作重新实现它。各映射及其像保持常秩的区域可以构成商向量丛；秩变化处应分层处理并重新验收，不能声称存在覆盖所有姿态的光滑固定秩丛。超声的中心与偏心切换恰是这样的秩变化例。

\section{内部解剖趋势、几何投影规则与示教手感的分工}
\label{sec:v43:anatomy-rules}
\subsection{血管和脏器共享接口，但不共享一条中心线}
腿部与躯干分别使用各自有效调和圆柱图；内部任务仍以部位、结构和进度为索引。血管先验可以在每个离散分支 $b$ 内写为
\begin{equation}
 \xi_b(s)=\bar\xi_b(s;c)+g_b(s),\quad
 g_b\sim\mathcal{GP}(0,K_{{\rm anat},b}),\qquad
 \Pr(b\mid\mathcal D)=\pi_b,
\end{equation}
再由 $P_{a,c,f}(\xi_b)$ 映射到物理空间。图接缝附近使用周期角核或连续 lift；各分支保留关联，不能将两个血管分支的高斯均值混成一条实际不存在的血管。若先验来自解剖知识而非已采集数据，也可以直接给概率分布或确定性范围，不要求先训练 GP。

脏器任务可由材料域中的隐式边界 $\phi_o(\xi)=0$、目标区域 $\{\phi_o\le0\}$ 和一族指定扫掠切面表示。扫描参数 $s$ 索引一个当前扫掠段，切面编号保留为离散任务标签。脏器不是一条长血管；使用相同图册、纤维和任务实现接口，并不要求把脏器改写成中心线。呼吸滑移或接触压迫导致的内部位移须保留独立先验和在线估计。

\subsection{投影规则保留完整内部几何，GP 只学习相对动作}
对一个内部目标 $p_{\rm int}(s)$，沿已定义的局部路径或法向场找到体表作用点 $p_{\rm surf}(s)$。投影必须在所用片区内唯一、连续且满足接触可行性；若不存在这样的射线交点，切换片区或报告该规则不可用。世界空间中的血管切向为
\begin{equation}
 t(s)=\frac{\partial_s P_{a,c,f}(\xi_b(s))}
 {\|\partial_s P_{a,c,f}(\xi_b(s))\|},
 \quad R_{\rm geom}(s)=\mathcal R\bigl(t(s),n_{\rm surf}(s),\mathcal T\bigr),
 \quad \rho_{\rm geom}=Q^\top R_{\rm geom}.
\end{equation}
$\mathcal R$ 是任务指定的成像平面与接触姿态规则，例如横断面、长轴面或允许的混合切面；该规则不必对所有 $t,n_{\rm surf}$ 都有精确同时满足的解，接近退化时需重新验收。不能先把内部路径的深度分量删除，再声称扭曲血管的完整姿态信息已经保留。内部切向、切面目的、体表法向和接触允许角共同决定几何默认动作。

示教提供的是相对于该默认动作的可测量偏好：体表材料位置偏置、相对旋转、在有力测量时的压力规律，以及推进节奏。定义
\begin{equation}
 p_{\rm CoP}^{\rm pref}=P_{a,c,f}(\xi_{\rm surf}^{\rm pref}),\qquad
 R^{\rm pref}=Q\rho_{\rm geom}\Upsilon,\quad\Upsilon\in SO(3).
\end{equation}
当 $Q'=QS$、$\rho_{\rm geom}'=S^\top\rho_{\rm geom}$ 时，$\Upsilon$ 不变。这个相对量很适合学习，它把几何导致的转动与专家额外动作分开。它不保证某个内部中心线先验正确，也不把几何转动自动认定为专家手感。

GP forcing 学习技能形状，解剖 GP 学习结构位置；两者使用同一种 GP 推断方法但数据语义和后验分开。动捕没有测力时，不报告已学到压力或物理阻抗。法向力参考若由同步示教得到，也要先投影到任务规定范围；机器人接触稳定性由已有接触控制与实际验收负责。

\subsection{技能容差不替代任务完成条件}
记 GP 名义技能为 $Y_{\rm pref}(s)$，可信参考度量中的偏离代价为
\begin{equation}
 J_{\rm style}(a)=\sum_j
 \big\|\operatorname{Log}_{Y_{\rm pref}(s_j)}Y(a,s_j)\big\|^2_{\Lambda_j}.
\end{equation}
$\Lambda_j$ 可来自预先定义的技能容差；若来自示教变化，须区分有效动作多样性、均值估计不确定性与测量噪声。后验方差大不表示任务允许随意偏离。真正必须保持的姿态、声窗、接触和相位条件仍写入 $\mathcal A_\tau,\mathcal Q_\tau$。IRD 可以在该集合中改变实现，不能通过丢弃一个任务特征来降低校准维数。

\section{唯一技能路线：GP forcing 与材料几何二阶生成}
\label{sec:v43-geometric-skill}

\subsection{取舍：在材料空间学习，相同技能允许不同物理实现}
本版固定一条技能路线：GP 学习相对于解剖几何基准的形状驱动项，固定增益的几何二阶系统生成名义技能。机器人世界笛卡尔位姿仅用于动作执行和对照，不作为迁移技能的主要学习坐标。原因不是笛卡尔 GP 不能拟合轨迹，而是人体的旋转、体型、姿态和体表曲率会被混入其训练目标；材料拉回则明确扣除这些已知变化。AL-DMP 已给出空间与速度分离的经典基础\cite{aldmp2016}；几何 DMP 已给出流形技能生成的基础\cite{gdmp2024}；GP 学习 DMP forcing 也已有超声接近动作先例\cite{gpdmp2026}。这些基础在本版中的用途是稳定表达和复现技能，不单独计为原创理论。

内部血管或脏器先验仍由解剖模型及在线估计提供，不能从专家探头动作反推为准确内部结构。给定任务及该先验，在体表材料子域 $\mathcal B\subset\mathcal M$ 上生成一条几何基准曲线 $\gamma(s)$，并由内部目标方向、合格切面与接触几何定义名义姿态 $\rho_{\rm geom}(s)$。这一步必须保留内部几何，单独将深度置零的体表投影一般不能保留切面朝向。专家示教学习的是相对于上述基准的偏置和姿态残余，而不是再次学习整个人体的世界位置。

令 $g_0$ 为预先固定的材料度量，$g_{\mathcal B}$ 为其在体表子域上的诱导度量，使用局部法邻域中的表达
\begin{equation}
 \xi_{\rm pref}(s)=\operatorname{Exp}_{\gamma(s)} b(s),\qquad
 \rho_{\rm pref}(s)=\rho_{\rm geom}(s,\xi_{\rm pref}(s))\Upsilon(s),
 \label{eq:v43-skill-state}
\end{equation}
其中 $b(s)\in T_{\gamma(s)}\mathcal B$ 是体表切向偏置，$\Upsilon(s)\in SO(3)$ 是几何基准之上的姿态残余。$\rho_{\rm geom}(s,\xi)$ 是定义在允许位置邻域上的几何姿态截面，而非只有中心线位置才能评价的姿态序列。$\operatorname{Exp}$ 只在验收过的单射半径内使用；大范围运动用重叠图和状态运输。腿部与躯干可以使用不同圆柱图，但可以共享此接口。共同接口不意味着腿部血管技能无需新数据就能外推到腹部器官。

曲线的进度取材料几何基准的归一化弧长，而非含纠偏往返的实际 TCP 弧长：
\begin{equation}
 \ell(u)=\int_0^u\|\gamma_0'(r)\|_{g_{\mathcal B}}\,dr,\qquad
 L_\gamma=\ell(1),\qquad s=\frac{\ell(u)}{L_\gamma},\qquad
 \dot s=\nu(t)\ge0.
 \label{eq:v43-material-progress}
\end{equation}
任务片段为 $s\in[0,1]$；先验的内部材料弧长也可以作为进度，但所有示教和生成必须使用相同定义。分叉是离散任务分支，不应被单个高斯均值连成穿过两支之间的虚假中心线。

\subsection{真实 connection 与姿态残余的 gauge 消去}
材料图册的被动换图为 $\xi'=\phi(\xi)$。沿 $\gamma$ 的切向量 $b$ 与其协变导数变换为
\begin{equation}
 b'=D\phi\,b,\qquad
 (D_sb)^i=\frac{db^i}{ds}
 +\Gamma^i{}_{jk}(\gamma)\gamma'^j b^k,
 \label{eq:v43-coordinate-connection}
\end{equation}
其中 $\Gamma$ 是 $g_{\mathcal B}$ 的 Levi--Civita connection；若使用三维材料表达，应将环境 connection 投影到体表切空间。对 $z_b=D_sb$ 再求 $D_sz_b$ 时也必须包含相同的 connection。调和坐标不是全局欧式坐标，不能将不同图中的分量普通二阶导数直接当成同一个加速度。本文的几何方程使用真实协变导数。附件 G-DMP 文献讨论了将协变导数近似为普通导数的实现\cite{gdmp2024}；本版不把该近似升级为任意换图下的严格不变性证明。

姿态纤维的标架改变为 $Q'=QS$、$\rho'=S^\top\rho$。连接一形式 $\widehat A=Q^\top dQ$ 满足
\begin{equation}
 \widehat A'=S^\top\widehat A S+S^\top dS.
 \label{eq:v43-gauge-connection}
\end{equation}
几何姿态同样运输为 $\rho_{\rm geom}'=S^\top\rho_{\rm geom}$，于是
\begin{equation}
 \Upsilon'=\rho_{\rm geom}'{}^\top\rho'
 =\rho_{\rm geom}^\top\rho=\Upsilon.
 \label{eq:v43-invariant-residual}
\end{equation}
这是选择相对姿态残余的重要好处：局部体表标架的任意旋转不会被学习成专家手感。它仍不是把整个物理姿态丛变成一个全局乘积；只是在存在有效几何参考截面的任务窗口内，对相对技能使用这一平凡化。

对实际身体映射后的基准姿态
\begin{equation}
 R_0(t)=Q_{a,c,f}(\xi(t))\rho_{\rm geom}(s(t),\xi(t)),\qquad
 R(t)=R_0(t)\Upsilon(t),
\end{equation}
设 $\omega_R=(R^\top\dot R)^\vee$ 和 $\omega_0=(R_0^\top\dot R_0)^\vee$，则示教的残余角速度为
\begin{equation}
 \omega_\Upsilon=\omega_R-\Upsilon^\top\omega_0.
 \label{eq:v43-human-angle-compensation}
\end{equation}
$\xi(t)$ 在示教时取实际 CoP 材料位置，在生成时取名义偏置位置。$\omega_0$ 包括沿材料路径、位置偏置变化以及实际人体运动引起的标架变化。平移示教也先按前文 $\dot\xi=J_P^{-1}(v_{\rm CoP}-b_p)$ 扣除人体运动，并使用 CoP、同步人体姿态和探头位姿。没有这个补偿，GP 会把被试移动误记为医生的主动纠偏。

\subsection{固定增益的空间二阶系统}
将偏置目标 $b_g(s)$ 定义为沿 $\gamma$ 平行运输的固定向量，即 $D_sb_g=0$；常用选择是 $b_g=0$。姿态残余目标 $\Upsilon_g$ 固定，常用选择为单位阵。令 $z_b=D_sb$，$\Upsilon'=\Upsilon\widehat\omega$，使用
\begin{align}
 D_sb&=z_b,\\
 D_sz_b&=-k_b(b-b_g)-d_bz_b+F_b(s),\\
 \frac{d\Upsilon}{ds}&=\Upsilon\widehat\omega,\\
 \frac{d\omega}{ds}
 &=k_\rho\operatorname{Log}(\Upsilon^\top\Upsilon_g)^\vee
   -d_\rho\omega+F_\rho(s).
 \label{eq:v43-space-generator}
\end{align}
所有 $k_b,d_b,k_\rho,d_\rho$ 是固定正数，不由 GP 学习。在局部线性模型下可取 $d_b=2\sqrt{k_b}$、$d_\rho=2\sqrt{k_\rho}$，但这不保证任意大角度姿态、饱和与物理执行下无超调。$SO(3)$ 使用与 $\|\omega\|^2$ 对应的双不变度量；沿同一曲线对其速度求协变导数时，body 表达中的 connection 项为 $\tfrac12[\omega,\omega]=0$，故此处的 $d\omega/ds$ 是精确的该速度协变加速度表达，而非忽略 connection 的一般近似。姿态 Log 使用 $\|\operatorname{Log}(\Upsilon^\top\Upsilon_g)^\vee\|<\pi-\epsilon$ 的有效窗口。

增益各向同性使被动正交标架变化直接保持方程。若需要矩阵增益，则 $K_b,D_b$ 必须是协变的正定张量，并随平行运输更新；不允许在每个图中重新指定相同数值矩阵而宣称同一物理技能。

时间执行直接乘进度速度：
\begin{align}
 \dot s&=\nu,\qquad D_tb=\nu z_b,\\
 D_tz_b&=\nu[-k_b(b-b_g)-d_bz_b+F_b(s)],\\
 \dot\Upsilon&=\nu\Upsilon\widehat\omega,\\
 \dot\omega&=\nu[k_\rho\operatorname{Log}(\Upsilon^\top\Upsilon_g)^\vee
                    -d_\rho\omega+F_\rho(s)].
 \label{eq:v43-time-generator}
\end{align}
对所有状态同时缩放是关键。$\nu=0$ 时材料参考与其内部形状状态全部暂停，避免只冻结 canonical phase、弹簧仍向终点拉动。实际人体变化仍通过 $P_{a,c,f}$ 和 $Q_{a,c,f}$ 运输该材料参考，所以暂停材料推进不等于将世界位姿钉死。接触导纳与有效观测下的局部纠偏仍继续执行。

空间曲线对任意单调时间重参数保持相同；其物理速度与加速度并不保持相同。例如
\begin{equation}
 D_t^2 b=\dot\nu z_b+\nu^2D_sz_b.
 \label{eq:v43-time-acceleration}
\end{equation}
因而 $\nu$ 和 $\dot\nu$ 仍需按实际加速度、接触与 IRD 能力限制平滑调度。任务反向重扫应切换到独立反向技能或重新初始化；本版不对 $\nu<0$ 宣称耗散稳定。

\subsection{forcing 示教标签：共同切空间与导数噪声}
平移向量在不同 $s$ 属于不同切空间。令 $\mathsf P_{s\leftarrow0}$ 是沿同一基准曲线的 Levi--Civita 平行运输，定义共同参考表达
\begin{equation}
 \bar b(s)=\mathsf P_{0\leftarrow s}b(s),\quad
 \bar z_b=\frac{d\bar b}{ds},\quad
 F_b(s)=\mathsf P_{s\leftarrow0}\bar F_b(s).
 \label{eq:v43-common-tangent}
\end{equation}
选取 $T_{\gamma(0)}\mathcal B$ 的正交基后，位置 forcing 标签是
\begin{equation}
 \boxed{\bar F_{b,d}(s)=\bar b_d''(s)
      +d_b\bar b_d'(s)+k_b[\bar b_d(s)-\bar b_g].}
 \label{eq:v43-position-label}
\end{equation}
这里的导数是完成材料拉回、平行运输及相位对齐之后的空间导数。$\dot s$ 太小的重复样本先合并，不在停顿期间用 $\dot b/\dot s$ 制造无限标签。

姿态标签由
\begin{equation}
 \boxed{F_{\rho,d}=\omega_d'
       -k_\rho\operatorname{Log}(\Upsilon_d^\top\Upsilon_g)^\vee
       +d_\rho\omega_d}
 \label{eq:v43-orientation-label}
\end{equation}
得到。若平滑示教使用 $r_d(s)=\operatorname{Log}(\Upsilon_g^\top\Upsilon_d(s))^\vee$，则
\begin{align}
 \omega_d&=J_r(r_d)r_d',\\
 \omega_d'&=J_r(r_d)r_d''+
              DJ_r(r_d)[r_d']r_d',\\
 J_r(r)&=I-\frac{1-\cos\theta}{\theta^2}\widehat r
              +\frac{\theta-\sin\theta}{\theta^3}\widehat r^2,
       \qquad\theta=\|r\|.
 \label{eq:v43-right-jacobian}
\end{align}
零角附近用解析极限。对原始欧拉角直接差分，不等于这个角速度与协变加速度。

GP 采用共同切空间中的输出 $(\bar F_b,F_\rho)$。旋转输出位于固定的 Lie algebra $\mathfrak{so}(3)\simeq\R^3$，再用左平移 $F_\rho\mapsto\Upsilon\widehat F_\rho\in T_\Upsilon SO(3)$。这是明确的 tangent trivialization，不是把处在不同姿态的切向量未经说明相加。位置运输与旋转 left trivialization 可组成一个低维矩阵值 GP；共用一种学习方法不要求将解剖位置先验和专家 forcing 混成同一个后验。

二阶差分容易放大动捕噪声。对 $\bar b$ 先拟合可二次均方微分的 GP 平滑器，记
\begin{equation}
 \mathcal L_b=\partial_s^2+d_b\partial_s+k_b.
\end{equation}
若平滑后的条件均值和协方差为 $m_b,K_b^{\rm post}$，则线性导数推断给出
\begin{align}
 m_{F_b}(s)&=\mathcal L_bm_b(s)-k_b\bar b_g,\\
 \operatorname{Cov}(F_b(s),F_b(s'))
 &=\mathcal L_{b,s}\mathcal L_{b,s'}K_b^{\rm post}(s,s').
 \label{eq:v43-derivative-noise}
\end{align}
这一标签噪声在邻近相位之间相关，训练 forcing GP 时应保留该协方差或明确给出其近似。Mat\'ern $5/2$ 核满足本段所需的二阶均方导数。姿态式~\eqref{eq:v43-right-jacobian} 是非线性变换，不能把其后验称为精确高斯；可用局部 Jacobian 传播或后验样本生成标签及噪声估计。

示教平滑、标签变换、再训练 forcing GP 是两阶段 plug-in／经验贝叶斯构造。下文 GP 后验条件于所采用的标签和噪声模型；除非另行建立联合概率模型，不宣称它是原始示教的精确联合后验。这不影响一致运输后的标架协变结论。

本版不使用 $(g-y_0)$ 的除法归一化。零起终点偏置和闭环姿态残余都可能有 $g=y_0$；forcing 标签仍由式~\eqref{eq:v43-position-label}--\eqref{eq:v43-orientation-label}直接定义。若希望无限进度下 forcing 渐消，令
\begin{equation}
 \mathcal K_F(s,s')=\chi(s)\chi(s')k_f((s,c),(s',c'))\,\Sigma_F,
 \qquad \chi(s)=e^{-\alpha s},\quad\alpha>0,
 \label{eq:v43-gated-kernel}
\end{equation}
并直接对完整 forcing 标签建模，不把接近零的标签再除以 $\chi$。$\Sigma_F\succeq0$ 可保留偏置与旋转驱动间的相关性。有限扫描末端 $s=1$ 的 $\chi$ 一般不为零，也不需要为零。

\subsection{可计算 GP 与 atlas/gauge 协变}
堆叠标签 $\mathbf f$，令 $K$ 为上述矩阵核的训练 Gram 矩阵，$\Sigma_\epsilon$ 为示教标签噪声协方差。预测为
\begin{align}
 m_F(x_*)&=m_0(x_*)+K_{*D}(K+\Sigma_\epsilon)^{-1}
                         [\mathbf f-\mathbf m_0],\\
 K_F^{\rm post}(x_*,x_*')&=K_{**'}-
     K_{*D}(K+\Sigma_\epsilon)^{-1}K_{D*'}.
 \label{eq:v43-forcing-gp}
\end{align}
生成器使用 $m_F$，不需要从预测分布随机采样实际探头动作。选择新的人体上下文 $c$ 时，训练数据不足覆盖的部分会保留不确定性；少量专家数据只能支持明确的低维技能类，不能保证任意器官泛化。

若共同参考正交基更换，输出向量为 $\mathbf f'=O\mathbf f$，均值和协方差同步变为 $m_F'=Om_F$、$K_F'=OK_FO^\top$。在原始相位相关局部标架中，一般运输矩阵为 $T(s)$，对应矩阵核必须满足
\begin{equation}
 \mathcal K_F'(s,s')=T(s)\mathcal K_F(s,s')T(s')^\top.
 \label{eq:v43-gauge-kernel}
\end{equation}
标签噪声和固定增益也作同一变换，故式~\eqref{eq:v43-forcing-gp}的预测与变换交换，生成的物理偏置保持相同。流形向量值 GP 文献已有相关的 gauge-independent 建模基础\cite{gpmanifold2021}；本版使用任务窗口内的显式运输即可，不额外引入大网络。运输沿指定曲线定义；不同绕行路径可能存在 holonomy，因此不声称平行运输与路径无关。

\subsection{名义生成器的能量性质与终点边界}
令 forcing 是预测均值，取能量
\begin{equation}
 \mathcal E=\tfrac12\|z_b\|^2+
 \tfrac{k_b}{2}\|b-b_g\|^2+
 \tfrac12\|\omega\|^2+
 \tfrac{k_\rho}{2}d_{SO(3)}(\Upsilon,\Upsilon_g)^2.
 \label{eq:v43-skill-energy}
\end{equation}
向量范数均为固定材料度量与旋转度量下的范数，距离仅在上述 Log 有效域内求导。$D_sb_g=0$ 和 connection 的度量相容性给出
\begin{equation}
 \mathcal E'=-d_b\|z_b\|^2-d_\rho\|\omega\|^2
               +\langle z_b,F_b\rangle+\omega^\top F_\rho.
 \label{eq:v43-energy-identity}
\end{equation}
设 $d_{\min}=\min(d_b,d_\rho)$，Young 不等式得到
\begin{equation}
 \mathcal E'\le -\tfrac{d_{\min}}2
          (\|z_b\|^2+\|\omega\|^2)
             +\frac{\|F_b\|^2+\|F_\rho\|^2}{2d_{\min}}.
 \label{eq:v43-energy-inequality}
\end{equation}

\begin{proposition}[名义稳定与有限段能量界]
\label{prop:v43-nominal-skill}
固定基准、固定目标、固定增益下，若轨迹保留在一个不接触姿态 cut locus 的紧致有效邻域内，则零 forcing 的平衡 $(b_g,0,\Upsilon_g,0)$ 在该局部域中渐近稳定。若 forcing 有界且趋零、平方可积，并且解、空间速度及其一阶导数均匀连续（例如相应一、二阶空间导数有界），则延拓到 $s\to\infty$ 的解趋向该平衡。在任意有限进度段 $[0,S]$，只需 forcing 有界，即有
\begin{equation}
 \mathcal E(S)\le\mathcal E(0)+
  \frac1{2d_{\min}}\int_0^S
     (\|F_b\|^2+\|F_\rho\|^2)\,ds.
 \label{eq:v43-finite-energy}
\end{equation}
\end{proposition}
\begin{proof}
式~\eqref{eq:v43-energy-identity}在 forcing 为零时只允许零速度留在 $\mathcal E'=0$ 的最大不变集；动力学随后要求偏置目标和姿态目标同时满足。有效域中唯一的势能临界点是上述平衡，故局部 LaSalle 论证成立。式~\eqref{eq:v43-energy-inequality}直接积分得到有限段能量界。对平方可积 forcing，该式给出速度平方可积；速度均匀连续使其趋零，速度导数均匀连续及速度极限存在进一步使速度导数趋零。代回动力学并使用 forcing 趋零，得到 $b-b_g\to0$ 和 $\operatorname{Log}(\Upsilon^\top\Upsilon_g)\to0$。所有论证均以有效域保持为前提；可用初始能量与 forcing 能量小于域边界上的最低势能来验收一个充分条件。
\end{proof}

时间能量变化为 $\dot{\mathcal E}=\nu\mathcal E'$。因此暂停不增加名义生成器能量。但若累计进度 $\int_0^\infty\nu\,dt$ 有限，尤其到 $s=1$ 就停，不能用 $s\to\infty$ 的结果宣称已经到达吸引子。有限扫描末端必须另外规定允许的偏置、姿态和速度容差；可采用经检查的末端条件化或短段接续。端点位姿或速度不兼容时，不能仅靠 forcing 渐消自动修复。基准目标更新、人体重注册和分叉切换同样应做状态运输与平滑接续，不能搬用固定基准的能量证明。

\subsection{与 IRD、反馈修复及失锁规则的完整接口}
技能生成器只输出 $\xi_{\rm pref},\rho_{\rm pref}$ 及其导数；它不取消允许集合 $\mathcal A_\tau$。物理候选仍为
\begin{equation}
 p_{\rm CoP}=P_{a,c,f}(\xi),\qquad
 R=Q_{a,c,f}(\xi)\rho,\qquad a\in\mathcal A_\tau.
\end{equation}
IRD 在这些允许实现中提出候选。可用材料偏置与姿态残余距离给出技能偏好，例如
\begin{equation}
 \mathcal L_{\rm skill}(a,s)=
 \lambda_b\|b(a)-b_{\rm pref}(s)\|^2+
 \lambda_\rho d_{SO(3)}(\Upsilon(a),\Upsilon_{\rm pref}(s))^2,
 \label{eq:v43-skill-preference}
\end{equation}
但同一个有效任务、允许偏离、真实观测质量、接触、输入能力和碰撞条件仍是硬验收。候选内逐一重算 $H,N,D_\star,D_Q$，探查数据仍由实际执行产生，反馈修复仍由 $q<1$ 的证书验收。因此更换技能编码不改变不可消除耦合商空间、最小校准方向及任务恢复的核心主张。

名义 forcing 的能量证明不等于患者端口无源、物理力稳定，也不证明 GP 方差构成可靠误差管。外层图像纠偏、接触导纳、QP 限幅和机器人延迟仍遵守各自证书。二阶参考生成与二阶图像纠偏作用在不同状态：前者表达希望的材料技能，后者修复实际观测误差；不把两套弹簧无说明地串在同一误差通道。图像失锁时，令 $\nu=0$ 冻结全部材料技能参考；图像纠偏状态制动，接触控制保持，未来误差管按缺少新观测的界扩张。恢复关联和声窗后才重新验收推进。

示教仅提供低维动作残余及可选推进偏好。合成内部血管可以训练解剖趋势、目标几何与仿真先验，但不能替代专家真实扫描动作作为“手感”证据。少量专家示教是否足够，应以留出人体的技能保持误差、暂停接续、有效声窗、接触变化和重新校准预算验证，而非规定一个未经实验支持的数据量。最强取舍对照是同一材料坐标、同一 IRD、同一观测、相同带宽及算力下的 GP 材料参考加二阶跟踪；若 forcing 编码没有额外收益，应承认二者的功能等价，不将符号丰富度当成新能力。

\section{局部响应模型：任务保持与不确定通道}
\label{sec:cal-local}

本节讨论一个明确且可检验的局部模型类。所有变量已经按物理运动度量与任务容差归一化，因而下文的欧氏范数、正交投影和奇异值具有相同的物理解释。固定允许的机器人实现后，省略其索引，令
\begin{equation}
  \dot y=Ju+w,\qquad
  J=H+\Theta N+E_{\mathrm{geo}},
  \label{eq:cal-model}
\end{equation}
其中 $y\in\mathbb R^m$ 是实际提取的有效任务特征，$u\in\mathbb R^n$ 是扣除人体运动并运输到材料标架后的实际实现运动。$H\in\mathbb R^{m\times n}$ 是已知名义响应，$N\in\mathbb R^{r\times n}$ 是已知的不确定物理通道，$\Theta\in\mathbb R^{m\times r}$ 是患者相关的未知响应。$E_{\mathrm{geo}}$ 汇集几何运输、模型遗漏和局部线性化的误差。扰动 $w$ 可以包括不能解释为命令响应的目标漂移、人体运动估计残差以及参考变化。

\begin{assumption}[局部模型、有效观测与物理执行]
\label{ass:cal-model}
在本次校准和验收的局部窗口内，$H$ 满行秩，$N$ 已知，且
\begin{equation}
  \|\Theta\|_2\le\delta,\quad\delta>0,\qquad
  \|E_{\mathrm{geo}}\|_2\le\eta .
\end{equation}
校准测量来自有效的目标特征；被实际执行的命令、而非理想规划命令，被用于构造数据。模型只在明确的控制与观测可信区间内使用。失锁或关联未确定的图像不产生新的有效任务响应测量。
\end{assumption}

$N$ 必须来自物理响应结构或经过验证的辨识结果。例如，特定接触变形通道可以由明确的接触模型给出。SMPL-X 对应关系和坐标不变性本身不会给出已知 $N$，也不会保证小的 $\eta$。若所有命令方向都任意未知，应取 $N=I_n$；此时不能预期以下结构产生省校准收益。首先在 $E_{\mathrm{geo}}=0$、无噪声且局部 $\Theta$ 固定的模型内证明精确结论，再给出带误差的验收条件。

\subsection{保持任务的物理实现}

名义任务保持的线性实现 $B\in\mathbb R^{n\times m}$ 满足
\begin{equation}
  HB=I_m,\qquad u=Bv ,
  \label{eq:cal-lift}
\end{equation}
其中 $v\in\mathbb R^m$ 是同一个材料任务命令。实际响应为
\begin{equation}
  JB=I_m+\Theta NB+E_{\mathrm{geo}}B .
  \label{eq:cal-exposure}
\end{equation}
下文精确谱、首次下界与必要充分方向数采用 $E_{\mathrm{geo}}=0$、无噪声的冻结局部模型；有余项时使用后文 $q$ 的验收证书。

因此，需要校准的是 $\Theta$ 在 $\operatorname{range}(NB)$ 上的作用，称 $NB$ 为该实现的不确定曝光。这里不要求未知矩阵 $\Theta$ 低秩；低维性来自可选择的实际机器人实现。

若精确接触条件要求 $Eu=0$，令 $Z_E$ 是 $\ker E$ 的正交基，在允许的命令空间内使用
\begin{equation}
  H_E=HZ_E,\qquad N_E=NZ_E .
  \label{eq:cal-contact-reduction}
\end{equation}
只有 $H_E$ 满行秩时，才能继续实现全部任务方向。以下代数结果可以对 $(H_E,N_E)$ 使用。实际探查所用输入来自运动测量；若只能获得速度接口指令，应将指令与实际运动的差异纳入独立执行误差界。速度、关节、碰撞和接触的不等式能力限制则需要另外验收；它们不能从右逆关系推得。

\section{任务实现如何减少必须校准的方向}
\label{sec:cal-realization}

\subsection{同时最小化所有奇异值的实现}

取 $Z_H$ 为 $\ker H$ 的正交基。由于 $H$ 满行秩，$HH^\dagger=I_m$，而任意右逆都有形式
\begin{equation}
  B=H^\dagger+Z_HK .
  \label{eq:cal-all-lifts}
\end{equation}
记
\begin{equation}
  A=NH^\dagger,\qquad C_0=NZ_H,\qquad
  \mathcal C=\operatorname{range}(C_0)=N\ker H .
\end{equation}
令 $P_{\mathcal C}=C_0C_0^\dagger$ 为到 $\mathcal C$ 的正交投影，并定义
\begin{equation}
  B_*=
  H^\dagger-Z_HC_0^\dagger A,\qquad
  D_*=
  NB_*=(I-P_{\mathcal C})A .
  \label{eq:cal-optimal-lift}
\end{equation}

\begin{theorem}[最小不确定曝光]
\label{thm:cal-singular-opt}
在无额外物理不等式约束的右逆集合 $HB=I_m$ 中，$B_*$ 同时最小化 $NB$ 的所有奇异值：
\begin{equation}
  \sigma_j(NB)\ge\sigma_j(D_*),
  \qquad j=1,\ldots,\min\{r,m\}.
  \label{eq:cal-singular-order}
\end{equation}
特别地，它同时最小化谱范数、任意奇异值尾和曝光的秩。
\end{theorem}

\begin{proof}
对任意 $B=H^\dagger+Z_HK$，
\begin{align}
  NB
  &= A+C_0K\\
  &= (I-P_{\mathcal C})A+
     \big(P_{\mathcal C}A+C_0K\big)\\
  &=D_*+E_B .
\end{align}
$D_*$ 的列属于 $\mathcal C^\perp$，$E_B$ 的列属于 $\mathcal C$，故 $D_*^\top E_B=0$。于是
\begin{equation}
  (NB)^\top NB
  =D_*^\top D_*+E_B^\top E_B
  \succeq D_*^\top D_* .
\end{equation}
由对称矩阵的特征值单调性得到~\eqref{eq:cal-singular-order}。选择 $K=-C_0^\dagger A$ 时 $E_B=0$，从而各个下界均达到。
\end{proof}

这个结论区分了两种改变：被动坐标变化不会创造新的曝光收益；主动改变物理命令实现，可以用 $N\ker H$ 中的额外自由度抵消不确定曝光，而同时保留 $HB=I_m$。

\subsection{最低精确校准维数}

\begin{theorem}[任务保持实现的结构秩]
\label{thm:cal-rank}
最低不确定曝光秩为
\begin{equation}
  \boxed{
  k_*=
  \min_{HB=I_m}\operatorname{rank}(NB)
  =m-\dim(H\ker N)
  =m+\operatorname{rank}N-
  \operatorname{rank}
  \begin{bmatrix}H\\N\end{bmatrix}.}
  \label{eq:cal-kstar}
\end{equation}
\end{theorem}

\begin{proof}
因为 $HB=I_m$，$B$ 是单射，而且 $H$ 在 $\operatorname{range}(B)$ 上是到 $\mathbb R^m$ 的同构。因此
\begin{align}
  \operatorname{rank}(NB)
  &=m-\dim\big(\operatorname{range}(B)\cap\ker N\big)\\
  &\ge m-\dim(H\ker N).
  \label{eq:cal-rank-lower}
\end{align}
另一方面，$D_*v=0$ 当且仅当 $NH^\dagger v\in N\ker H$。这等价于存在 $z\in\ker H$ 使
\begin{equation}
  N(H^\dagger v-z)=0,\qquad
  H(H^\dagger v-z)=v ,
\end{equation}
即 $v\in H\ker N$。因而
\begin{equation}
  \ker D_*=H\ker N,\qquad
  \operatorname{rank}D_*=m-\dim(H\ker N),
\end{equation}
下界由 $B_*$ 达到。最后，对从命令空间到联合输出 $(Hu,Nu)$ 的线性映射，先取 $Nu$ 的像，再加入 $Nu=0$ 时的 $Hu$，得到
\begin{equation}
  \operatorname{rank}
  \begin{bmatrix}H\\N\end{bmatrix}
  =\operatorname{rank}N+\dim(H\ker N).
\end{equation}
\end{proof}

当 $N=I_n$ 时，$H\ker N=\{0\}$，故 $k_*=m$。因此，结构收益依赖实际不确定通道与任务通道之间的关系，而不能由改变参数化自动获得。$B_*$ 也可能增加所需物理运动幅度；受约束的实际问题必须保留这一代价。

\section{允许重新选择任意物理动作的首次任务下界}
\label{sec:cal-minimax}

这一节不预先锁定 $u=Bv$。校准以后，控制器可以根据数据重新选择右逆，甚至使用任意非线性的静态规则 $u(v)$。因此，需要剔除任务零空间提供的所有绕开机会，不能把固定 $B$ 的谱尾直接当作全体实现的任务下界。

\subsection{校准信息与正确的剩余空间}

一条无噪声局部校准记录给出实际响应增量 $\Delta y_i$（或对冻结 Jacobian 积分的局部响应）。输入 $u_i$ 对应相同时间尺度的实际积分运动，因此
\begin{equation}
  r_i=\Delta y_i-Hu_i=\Theta x_i,\qquad x_i=Nu_i .
  \label{eq:cal-probe-oracle}
\end{equation}
记 $Q=\operatorname{span}\{x_1,\ldots,x_s\}$。这里计数的是独立激励方向 $\dim Q$，不是图像帧数、探查持续时间或动作片段数。长片段可以提供多个独立激励。

考虑所有校准残差均为零的历史。数据一致模型满足
\begin{equation}
  \Theta Q=0,\qquad \|\Theta\|_2\le\delta .
  \label{eq:cal-zero-history}
\end{equation}
定义
\begin{equation}
  D_Q=P_{\mathcal C+Q}^{\perp}NH^\dagger,
  \qquad \mathcal C=N\ker H .
  \label{eq:cal-DQ}
\end{equation}
正确剩余矩阵是 $D_Q$，一般不是 $(I-P_Q)D_*$。实际控制器还能够通过 $\ker H$ 中的命令，把未校准曝光移动到已校准空间。

\begin{lemma}[给定名义任务输出的最小剩余曝光]
\label{lem:cal-DQ}
对任意 $z\in\mathbb R^m$，
\begin{equation}
  \min_{Hu=z}\|(I-P_Q)Nu\|_2
  =\|D_Qz\|_2 .
  \label{eq:cal-exposure-z}
\end{equation}
\end{lemma}

\begin{proof}
所有满足 $Hu=z$ 的命令为 $u=H^\dagger z+Z_H\ell$。记 $P_Q^\perp=I-P_Q$，则
\begin{align}
  \min_{Hu=z}\|P_Q^\perp Nu\|_2
  &=
  \min_\ell
  \|P_Q^\perp Az+P_Q^\perp C_0\ell\|_2\\
  &=
  \operatorname{dist}
  \big(P_Q^\perp Az,P_Q^\perp\mathcal C\big)\\
  &=\|P_{\mathcal C+Q}^\perp Az\|_2 .
\end{align}
最后一步因为 $Q^\perp$ 内再剔除 $P_Q^\perp\mathcal C$，等价于在原空间同时剔除 $\mathcal C+Q$。
\end{proof}

\subsection{非线性共同动作也不能绕开的下界}

定义零校准历史之后首次重新实现任务的最小最大误差
\begin{equation}
  \mathcal L(Q)=
  \inf_{u(\cdot)}
  \sup_{\|v\|_2=1}
  \sup_{\substack{\Theta Q=0\\\|\Theta\|_2\le\delta}}
  \|(H+\Theta N)u(v)-v\|_2 .
  \label{eq:cal-first-loss}
\end{equation}
这里允许任意静态非线性命令规则 $u(v)$，且暂时放松所有物理输入限制。因而获得的下界对有约束系统仍然有效；匹配的上界只有在相应命令确实可执行时才能使用。

\begin{theorem}[首次任务响应的精确条件最小最大误差]
\label{thm:cal-first-loss}
在~\eqref{eq:cal-zero-history} 的数据一致模型中，
\begin{equation}
  \boxed{\mathcal L(Q)=
  \min\{1,\delta\|D_Q\|_2\}.}
  \label{eq:cal-first-exact}
\end{equation}
\end{theorem}

\begin{proof}
对固定物理命令 $u$，令 $z=Hu$。由于 $\Theta Q=0$，$\Theta Nu=\Theta P_Q^\perp Nu$。范数球内的矩阵可以把 $P_Q^\perp Nu$ 映射到任意半径为 $\delta\|P_Q^\perp Nu\|_2$ 的任务输出向量。因此
\begin{equation}
  \sup_{\Theta Q=0,\;\|\Theta\|_2\le\delta}
  \|(H+\Theta N)u-v\|_2
  =
  \|z-v\|_2+\delta\|P_Q^\perp Nu\|_2 .
  \label{eq:cal-task-ball}
\end{equation}
利用引理~\ref{lem:cal-DQ}，逐个任务命令的最优值为
\begin{equation}
  \inf_z\big(\|z-v\|_2+\delta\|D_Qz\|_2\big).
  \label{eq:cal-task-prox}
\end{equation}
令 $v_*$ 为 $D_Q$ 的最大右奇异向量，$w_*$ 为相应左奇异向量，$s=\|D_Q\|_2$。当 $s>0$ 时，对任意 $z$，
\begin{align}
  \|z-v_*\|_2+\delta\|D_Qz\|_2
  &\ge
  |v_*^\top z-1|
  +\delta|w_*^\top D_Qz|\\
  &=
  |t-1|+\delta s|t|,
  \qquad t=v_*^\top z\\
  &\ge\min\{1,\delta s\}.
\end{align}
若 $\delta s\le1$，选择满足 $HB_Q=I_m$ 且 $P_Q^\perp NB_Q=D_Q$ 的最小曝光右逆，例如
\begin{equation}
  B_Q=
  H^\dagger-
  Z_H(P_Q^\perp NZ_H)^\dagger
  P_Q^\perp NH^\dagger ,
  \label{eq:cal-BQ}
\end{equation}
并令 $u(v)=B_Qv$，其最大误差不超过 $\delta s$。若 $\delta s\ge1$，选择零增量命令 $u(v)=0$，其误差为 $1$。当 $s=0$，$B_Q$ 实现零误差。故上下界匹配。
\end{proof}

可以直接构造一对限制任何物理动作的模型。选 $w_*\in(\mathcal C+Q)^\perp$ 满足 $D_Q^\top w_*=sv_*$。因为 $w_*^\top N$ 在 $\ker H$ 上为零，
\begin{equation}
  w_*^\top N
  =w_*^\top NH^\dagger H
  =sv_*^\top H .
  \label{eq:cal-adversary-factor}
\end{equation}
令
\begin{equation}
  \Theta_\pm=\pm\delta v_*w_*^\top .
\end{equation}
两模型都满足校准数据，但
\begin{equation}
  J_\pm=
  (I_m\pm\delta s v_*v_*^\top)H .
  \label{eq:cal-adversary-J}
\end{equation}
对任意物理命令 $u$，任务 $v_*$ 的误差投影都满足
\begin{equation}
  \max_{\pm}
  \left|
  (1\pm\delta s)v_*^\top Hu-1
  \right|
  \ge\min\{1,\delta s\}.
  \label{eq:cal-anyphysical-lower}
\end{equation}
这个构造同时约束全部任务零空间动作，消除了预先固定 $B$ 所产生的漏洞。

\subsection{最优探查方向与谱尾}

\begin{theorem}[最优独立校准方向的条件谱尾]
\label{thm:cal-width}
在无物理探查限制的零校准历史模型内，
\begin{equation}
  \min_{\dim Q\le k}\|D_Q\|_2
  =\sigma_{k+1}(D_*),
  \label{eq:cal-quotient-width}
\end{equation}
其中超过矩阵秩的奇异值取为零。因此
\begin{equation}
  \boxed{
  \min_{\dim Q\le k}\mathcal L(Q)
  =\min\{1,\delta\sigma_{k+1}(D_*)\}.}
  \label{eq:cal-minimax-width}
\end{equation}
\end{theorem}

\begin{proof}
因为 $D_*$ 的列在 $\mathcal C^\perp$ 中，
\begin{equation}
  D_Q=P_{\mathcal C+Q}^\perp D_* .
\end{equation}
投影掉 $\mathcal C$ 后，$Q$ 只能提供至多 $k$ 维的额外空间 $P_{\mathcal C}^\perp Q$。因此，左投影消除 $D_*$ 的至多 $k$ 个方向后，谱残差至少为 $\sigma_{k+1}(D_*)$。选择 $Q$ 为 $D_*$ 的前 $k$ 个左奇异向量张成的空间可以达到该下界。这些向量位于 $\operatorname{range}(N)$ 中，故存在物理命令满足 $Nu_i=q_i$。定理~\ref{thm:cal-first-loss} 完成最后一步。
\end{proof}

若验收要求首次任务误差不超过 $\epsilon<1$，零历史模型中最低方向数为
\begin{equation}
  k_\epsilon=
  \#\{j:\delta\sigma_j(D_*)>\epsilon\}.
  \label{eq:cal-epsilon-dimension}
\end{equation}
边界等号是否接受需与实验验收规则一致。

\begin{remark}[自适应探查、非零历史与闭环边界]
\label{rem:cal-width-scope}
对任何确定性自适应校准算法，或随机种子固定后的逐历史确定性保证，考虑其全部残差为零的历史。随机算法的期望风险或高概率复杂度不由此自动推出。最终 $Q$ 的维数不超过独立激励次数，而上述对手在该算法已经执行的全部探查上均返回零。因此~\eqref{eq:cal-minimax-width} 给出原始全模型类的必要最坏情况首次任务误差下界。

条件等式及其匹配上界仅针对零历史。非零观测会改变已知响应，并可能使数据一致不确定集合变为非对称或椭球集合，不能继续把原始 $D_*$ 的谱尾当作所有历史的精确上界。此外，首次任务响应下界不等于任意最终闭环恢复性能的下界；容许等待、继续适应或接受暂态误差的控制器解决的是另一个任务。
\end{remark}

\section{精确单位任务响应的校准复杂度}
\label{sec:cal-exact-complexity}

前一节已经给出零历史的精确近似响应结果。下面证明在完整无噪声模型类内，一个限定明确的任务保证具有匹配的校准复杂度。

\begin{theorem}[精确任务重新实现的必要与充分方向数]
\label{thm:cal-exact-task}
假定 $\delta\|D_*\|_2<1$。考虑校准后对一个非空任务命令邻域实现精确单位响应
\begin{equation}
  Ju(v)=v
  \label{eq:cal-unit-task}
\end{equation}
的规则，并要求该保证对所有 $\|\Theta\|_2\le\delta$ 成立。无约束模型类内，最坏情况下必须且只需 $k_*=\operatorname{rank}(D_*)$ 个独立校准方向。

物理限制不会削弱必要下界；充分性还要求相应探查和修复后的任务命令位于实际可执行可信区间内。
\end{theorem}

\begin{proof}
必要性可由定理~\ref{thm:cal-width} 在零误差下直接得到，也可用维数证明。取有效真实模型 $\Theta=0$。全部校准残差为零，任何满足 $\Delta Q=0$ 的小范数矩阵均与数据一致。若对全部这些模型都满足~\eqref{eq:cal-unit-task}，必须有
\begin{equation}
  Hu(v)=v,\qquad Nu(v)\in Q .
\end{equation}
因此能实现的任务命令属于 $H(N^{-1}Q)$，而
\begin{equation}
  \dim H(N^{-1}Q)
  \le\dim(H\ker N)+\dim Q .
\end{equation}
覆盖一个 $m$ 维邻域要求 $\dim Q\ge m-\dim(H\ker N)=k_*$。该证明允许任意静态非线性 $u(v)$，并适用于自适应算法的零历史。

充分性选用 $B_*$，令 $Q=\operatorname{range}(D_*)$，并校准 $\Theta$ 在 $Q$ 上的作用。由数据可完整知道
\begin{equation}
  G=JB_*=I_m+\Theta D_* .
\end{equation}
因为 $\|\Theta D_*\|_2\le\delta\|D_*\|_2<1$，$G$ 可逆。修复实现
\begin{equation}
  C_{\mathrm{exact}}=B_*G^{-1}
  \label{eq:cal-exact-repair}
\end{equation}
满足 $JC_{\mathrm{exact}}=I_m$。只需 $k_*$ 个独立校准方向即可识别 $G$，无需识别整个 $\Theta$。
\end{proof}

修复后的实际实现运动代价不能省略。例如
\begin{equation}
  \|C_{\mathrm{exact}}\|_2
  \le
  \frac{\|B_*\|_2}{1-\delta\|D_*\|_2}.
  \label{eq:cal-exact-command-bound}
\end{equation}
当分母小或 $B_*$ 运动幅度大时，校准方向减少仍可能以不可执行的命令为代价。

\subsection{零偏离探查为何不足}

\begin{proposition}[名义任务中性探查的最坏历史限制]
\label{prop:cal-neutral}
若所有探查均满足 $Hu_i=0$，则在零校准历史上 $Q\subseteq\mathcal C=N\ker H$，因而
\begin{equation}
  D_Q=D_*,
  \qquad
  \mathcal L(Q)=\min\{1,\delta\|D_*\|_2\}.
  \label{eq:cal-neutral-no-width}
\end{equation}
无论增加多少这种探查，都无法缩减该历史上的不可避免首次任务曝光。
\end{proposition}

\begin{proof}
$Hu_i=0$ 意味着 $u_i\in\ker H$，所以 $Nu_i\in\mathcal C$。于是 $\mathcal C+Q=\mathcal C$，直接得到结论。
\end{proof}

这使允许的有限任务偏离获得明确职责：它提供激励 $\mathcal C^\perp$ 中任务耦合通道的权限。该结论是最坏历史限制，而不是断言名义任务中性探查在每个实际患者上都无用。非零观测可能揭示能够用于新实现的任务零空间响应；仍须重新计算数据一致模型和实际能力。

\subsection{所有观测历史的一个充分上界}

取 $Q$ 为 $D_*$ 的前 $k$ 个左奇异方向，在任意无噪声历史上定义
\begin{equation}
  \widehat\Theta=\Theta P_Q,\qquad
  \widehat G=I_m+\widehat\Theta D_*,
  \qquad C_k=B_*\widehat G^{-1}.
\end{equation}
若 $\delta\|P_QD_*\|_2<1$，则
\begin{align}
  JC_k-I_m
  &=\Theta(I-P_Q)D_*\widehat G^{-1},\\
  \|JC_k-I_m\|_2
  &\le
  \delta\sigma_{k+1}(D_*)\,
  \|\widehat G^{-1}\|_2\\
  &\le
  \frac{\delta\sigma_{k+1}(D_*)}
  {1-\delta\|P_QD_*\|_2}.
  \label{eq:cal-nonzero-upper}
\end{align}
这是所有历史上的充分上界。它与零历史的必要下界通常存在间隙；不能将二者写成一个未经证明的精确样本复杂度。

\section{有噪声的实际探查与凸响应修复}
\label{sec:cal-noisy}

\subsection{使用实际执行量构造估计}

将经过物理验收并实际执行的探查命令堆叠为 $U$，令
\begin{equation}
  X=NU,\qquad
  R_{\mathrm{data}}=\Delta Y-HU=\Theta X+W,
  \qquad\|W\|_2\le\varepsilon_{\rm data} .
  \label{eq:cal-data}
\end{equation}
$\varepsilon_{\rm data}$ 属于具体探查输入 $U$ 和局部时间窗口；改变幅度、探查库或窗口后应重算该界。$W$ 包含测量噪声、重复探查间的未抵消漂移以及 $E_{\mathrm{geo}}U$。如果两次探查共享相同初始状态或有可信的基线消除，可降低漂移项；不能默认它自动消失。由此定义
\begin{equation}
  \widehat\Theta=R_{\mathrm{data}}X^\dagger,
  \qquad
  \widehat J=H+\widehat\Theta N,
  \qquad P_X=XX^\dagger .
  \label{eq:cal-estimate}
\end{equation}
因为
\begin{equation}
  \widehat\Theta
  =\Theta P_X+WX^\dagger,
\end{equation}
对任意修复实现 $C$，
\begin{align}
  (J-\widehat J)C
  &=
  \Theta(I-P_X)NC
  -WX^\dagger NC+E_{\mathrm{geo}}C,\\
  \|(J-\widehat J)C\|_2
  &\le
  \underbrace{
  \delta\|(I-P_X)NC\|_2}_{\text{未激励方向}}
  +
  \underbrace{
  \varepsilon_{\rm data}\|X^\dagger NC\|_2}_{\text{估计误差放大}}
  +
  \underbrace{
  \eta\|C\|_2}_{\text{结构外误差}} .
  \label{eq:cal-q}
\end{align}
记右侧为 $q(C,X)$。它同时检查是否激励了正确方向，以及激励是否具有足够的数值条件。单纯提高 $\operatorname{rank}(X)$ 不保证少量低幅度探查在噪声中有效。

\subsection{固定探查数据后的可计算修复}

设一个经实际运动学、碰撞和接触检查得到的输入内近似为
\begin{equation}
  \mathcal U=
  \{u:a_i^\top u\le b_i,\ i=1,\ldots,p\}.
\end{equation}
局部误差可信区间为 $\|e\|_2\le r_{\mathrm{task}}$，二阶纠偏的预定可信速度半径为 $r_\zeta$（下文由增广状态界验收），已知前馈命令为 $u_{\mathrm{ff}}$。固定实际探查数据后可解
\begin{equation}
\begin{aligned}
  \min_C\quad&
  \|C-C_{\mathrm{transported}}\|_{W_C}^2+
  \lambda\|C\|_F^2\\
  \text{s.t.}\quad&
  \widehat JC=I_m,\\
  &q(C,X)\le q_{\mathrm{goal}}<1,\\
  &EC=0,\\
  &a_i^\top u_{\mathrm{ff}}
   +r_\zeta\|a_i^\top C\|_2
   \le b_i,\quad i=1,\ldots,p.
\end{aligned}
\label{eq:cal-convex-repair}
\end{equation}
其中 $W_C$ 为固定正定代价度量，$EC=0$ 表达要求保留的线性接触约束。该问题是凸的：等式为仿射约束，$q$ 是线性矩阵映射的谱范数之和，能力条件是行范数约束。最后一组不等式保证所有可信纠偏速度状态下
\begin{equation}
  u=u_{\mathrm{ff}}+C\zeta\in\mathcal U\quad\text{当 }\|\zeta\|_2\le r_\zeta .
\end{equation}

若模型与物理检查使可行集为空，应增加有依据的探查、改变允许实现或报告该保证不可实现，而不是放宽未验证的误差半径。dIRD 的梯度可以用于提出候选实现，但其密度或可达性代理值不代替 $\mathcal U$ 的真实验收。

\section{经校准的二阶图像纠偏、能力验收与滚动未来管}
\label{sec:cal-recovery}

\subsection{参考生成与实际反馈是两个不同对象}

前文的 GP--AL--几何 DMP 生成器描述材料技能的名义形状、姿态偏好和推进节奏；它的内部状态不是实际超声特征的观测值。令 $y_r(t)$ 为已接受实现对应的任务特征参考，$y(t)$ 为当前可信图像及相关传感器提供的 $m$ 维特征，定义 $e=y-y_r$。技能输出可以包含更多自由度，但不能把不可观测的技能变量加入 $e$ 后声称实现了反馈。

在共同的物理白化、平行运输任务标架中，实际图像误差测量为 $\widehat e=e+n$。推荐控制器只增加一个任务纠偏速度状态 $\zeta\in\R^m$：
\begin{equation}
 \boxed{u=u_{\mathrm{ff}}+C\zeta,\qquad
 \dot\zeta=-\omega^2(e+n)-2\omega\zeta,\quad \omega>0.}
 \label{eq:cal-servo-second-order}
\end{equation}
它替换原来的静态误差到速度映射。$\zeta$ 是虚拟纠偏速度，$\omega$ 是固定设计带宽；这不是对患者施加的弹簧力。GP 学习参考生成器的形状驱动项；实际图像反馈的 $\omega$ 不由普通示教轨迹辨识，避免将两套弹簧重复作用于同一误差。

由修复约束 $\widehat JC=I_m$ 与式~\eqref{eq:cal-q}，定义
\begin{equation}
 \Delta=(J-\widehat J)C,\qquad
 \|\Delta(t)\|_2\le q<1,
 \qquad
 w=Ju_{\mathrm{ff}}-\dot y_r+w_{\mathrm{ext}}.
 \label{eq:cal-feedback-residual}
\end{equation}
在有效局部窗口内，实际误差动力学是
\begin{equation}
 \boxed{\dot e=(I_m+\Delta)\zeta+w,\qquad
 \dot\zeta=-\omega^2(e+n)-2\omega\zeta.}
 \label{eq:cal-error-dynamics}
\end{equation}
若采用 $u_{\mathrm{ff}}=C\dot y_r$，仍有 $w=\Delta\dot y_r+w_{\mathrm{ext}}$；不能默认名义前馈完全抵消未知响应。身体运动补偿误差、目标漂移、图册运输误差和实际执行偏差应合并进入经过验证的 $w$ 或 $n$。

在 $\Delta=0,w=n=0$ 时，$\dot e=\zeta$，所以
\begin{equation}
 \ddot e+2\omega\dot e+\omega^2e=0,
 \quad
 e(t)=\bigl[e(0)+(\zeta(0)+\omega e(0))t\bigr]e^{-\omega t}.
 \label{eq:cal-nominal-second-order}
\end{equation}
名义临界阻尼意味着这种局部模型的双实根，不意味着任意初始速度、执行延迟或耦合误差下所有分量都永不超调。无需对噪声图像误差做数值微分。

\subsection{从实际探查证书到二阶指数恢复}

\begin{theorem}[时变耦合与测量噪声下的二阶恢复]
\label{thm:cal-recovery}
假定在同一有效任务标架内 $\omega$ 固定、$\widehat JC=I_m$、$\|\Delta(t)\|_2\le q<1$，且
$\|w(t)\|_2\le w_{\max}$、$\|n(t)\|_2\le n_{\max}$。定义
\begin{equation}
 v=\zeta/\omega,\quad r=e+v,\quad
 \chi=\begin{bmatrix}r\\v\end{bmatrix},\quad
 V_s=\tfrac12\|\chi\|_2^2,\quad S=\sqrt{2V_s},
 \quad \lambda_s=\frac{\omega(1-q)}2.
 \label{eq:cal-servo-energy}
\end{equation}
令 $b_{\rm drv}=w_{\max}+\sqrt2\,\omega n_{\max}$，则
\begin{equation}
 \begin{aligned}
 S(t)&\le e^{-\lambda_st}S(0)
       +\frac{b_{\rm drv}}{\lambda_s}(1-e^{-\lambda_st}),\\
 \|e(t)\|_2&\le\sqrt2S(t),\qquad
 \|\zeta(t)\|_2\le\omega S(t).
 \end{aligned}
 \label{eq:cal-continuous-tube}
\end{equation}
此结论不要求 $\Delta$ 对称，也不要求其随时间恒定。
\end{theorem}

\begin{proof}
由式~\eqref{eq:cal-error-dynamics}直接变换得
\begin{equation}
 \dot r=-\omega r+\omega\Delta v+w-\omega n,
 \qquad
 \dot v=-\omega r-\omega v-\omega n.
 \label{eq:cal-servo-rv}
\end{equation}
于是
\begin{align}
 \dot V_s
 &=-\omega(\|r\|_2^2+\|v\|_2^2)
   +\omega r^\top(\Delta-I_m)v
   +r^\top w-\omega(r+v)^\top n\\
 &\le-\omega(\|r\|_2^2+\|v\|_2^2)
   +\omega(1+q)\|r\|_2\|v\|_2
   +S\bigl(\|w\|_2+\sqrt2\omega\|n\|_2\bigr)\\
 &\le-\omega(1-q)V_s+Sb_{\rm drv}.
 \label{eq:cal-servo-lyapunov}
\end{align}
最后一步使用 $2\|r\|_2\|v\|_2\le\|r\|_2^2+\|v\|_2^2$。在 $S>0$ 时除以 $S$，零点使用上 Dini 导数，得到
\begin{equation}
 D^+S\le-\lambda_sS+b_{\rm drv}.
\end{equation}
比较积分即得~\eqref{eq:cal-continuous-tube}。因为 $e=r-v$、$\zeta=\omega v$，得到误差与速度界。
\end{proof}

任务的最低恢复率 $\lambda_{\mathrm{task}}$ 与容许稳态误差 $\epsilon_{\mathrm{task}}$ 可以直接验收为
\begin{equation}
 \lambda_s\ge\lambda_{\mathrm{task}},\qquad
 \frac{\sqrt2b_{\rm drv}}{\lambda_s}\le\epsilon_{\mathrm{task}}.
 \label{eq:cal-task-rate}
\end{equation}
因此 $q$ 不只是响应模型的预测误差；它决定恢复速度、残差管与所需执行能力。提高 $\omega$ 会加快名义恢复，也会增加噪声驱动和速度需求。GP 后验方差描述技能或解剖先验的不确定性；$q(C,X)$ 描述探查后未知反馈响应的可验收误差。这两者不能互相替代。

\subsection{速度状态域替换旧静态能力条件}

取可信状态域 $S\le R_s$。若
\begin{equation}
 \sqrt2R_s\le r_{\mathrm{task}},\qquad
 b_{\rm drv}\le\lambda_sR_s,
 \label{eq:cal-servo-invariant}
\end{equation}
且域内观测、模型和接触均有效，则式~\eqref{eq:cal-continuous-tube}保证该状态域前向不变。定义 $r_\zeta=\omega R_s$，式~\eqref{eq:cal-convex-repair}的输入验收必须使用
\begin{equation}
 \boxed{a_i^\top u_{\mathrm{ff}}
       +r_\zeta\|a_i^\top C\|_2\le b_i,
       \quad i=1,\ldots,p.}
 \label{eq:cal-servo-capability}
\end{equation}
因为 $\|\zeta\|_2\le r_\zeta$，这保证所有该速度域中的命令都在真实输入内近似 $\mathcal U$ 中。固定 $r_\zeta$ 与探查数据时，能力条件仍为凸行范数约束；若将带宽、状态半径、实现与 $C$ 同时作为变量，则不能沿用原凸性结论。接触约束 $EC=0$ 使 $EC\zeta=0$，保留前文定义的线性接触相容条件，但不自动保证非线性声窗或患者端口无源。

若 QP、内环或限幅导致实际输入为 $u_{\mathrm{act}}=u+\delta u$，将 $J\delta u$ 计入 $w$ 并独立验证其界；若界不成立，应重规划或暂停。这只修正任务误差管。若硬约束作用于实际输入，且 $\delta u\in\mathcal D_u$，还必须使用紧化能力条件
\begin{equation}
 a_i^\top u_{\mathrm{ff}}+r_\zeta\|a_i^\top C\|_2
 +\sup_{d\in\mathcal D_u}a_i^\top d\le b_i.
\end{equation}
等价地，前文 $\mathcal U$ 应预先定义为按 $\mathcal D_u$ 紧化后的命令内近似。即使 $\zeta$ 连续，突变的 $C$ 或 $u_{\mathrm{ff}}$ 仍可使物理命令跳变，因此实现切换要连续接续并重新验收。时变正交标架应对 $e,\zeta$ 使用同一个度量兼容导数，或将二者平行运输到共同标架；其反对称连接项不改变二次能量。非正交坐标和时变度量不能忽略连接与度量导数。

\subsection{采样证书与围绕未来名义计划的集合管}

连续时间定理不覆盖任意数字实现。以显式 Euler 为例，在冻结一步模型中令 $a=h\omega$，有
\begin{equation}
 \chi_{j+1}=T_{\Delta_j}\chi_j+\eta_j,\qquad
 T_\Delta=
 \begin{bmatrix}(1-a)I_m&a\Delta\\-aI_m&(1-a)I_m\end{bmatrix}.
 \label{eq:cal-discrete-general}
\end{equation}
其中 $\eta_j$ 包括 $h[w_j-\omega n_j;-\omega n_j]$ 与有界采样、保持和模型变化余项。对任意采用的实际一步矩阵 $T_A$，可先检查
\begin{equation}
 T_A^\top P T_A\preceq\rho^2P,
 \quad P\succ0,\quad\rho<1
 \label{eq:cal-discrete-matrix-certificate}
\end{equation}
是否在全部允许模型上成立；只检查一个名义矩阵不足以认证未知模型。对于式~\eqref{eq:cal-discrete-general}及 $0<a<1$，两行分别取范数给出无需枚举 $\Delta$ 的充分界：
\begin{equation}
 \rho_h^2=\lambda_{\max}(G_h),\qquad
 G_h=\begin{bmatrix}
 (1-a)^2+a^2&a(1-a)(1+q)\\
 a(1-a)(1+q)&(1-a)^2+a^2q^2
 \end{bmatrix}.
 \label{eq:cal-discrete-rho}
\end{equation}
确实，$\|\chi_{j+1}-\eta_j\|_2^2$ 不超过
$[\|r_j\|_2,\|v_j\|_2]G_h[\|r_j\|_2,\|v_j\|_2]^\top$。并且
\begin{align}
 (I_2-G_h)_{11}&=2a(1-a)>0,\\
 \det(I_2-G_h)&=a^2(1-a)(1-q)
                   [3+q-a(1-q)]>0.
\end{align}
所以 $q<1,0<h\omega<1$ 时 $\rho_h<1$。这是冻结模型下的二阶证书，不能继续使用一阶公式 $1-h\kappa(1-q)$。实际采样内变化与图像延迟还必须进入 $\eta_j$ 的有效界。

为区分恢复包络与预测不确定性，令 $\bar\chi_{j+1}=T_0\bar\chi_j$ 为当前接受的名义未来，预测集合为
$\mathcal T_j^\chi=\bar\chi_j\oplus\mathbb B_2(d_j)$。由于
$T_{\Delta_j}-T_0$ 仅含右上块 $a\Delta_j$，若 $\|\eta_j\|_2\le\beta_j$，则
\begin{equation}
 \boxed{d_{j+1}=\rho_{h,j}d_j
             +h\omega q_j\|\bar v_j\|_2+\beta_j.}
 \label{eq:cal-tube-recursion}
\end{equation}
若参考生成器或已知外生驱动也进入名义预测，只需同步从 $\eta_j$ 中扣除同一已知项；不得重复计入。记 $b_j=h\omega q_j\|\bar v_j\|_2+\beta_j$，展开递推得
\begin{equation}
 d_{t+k}\le
 \Bigl(\prod_{j=t}^{t+k-1}\rho_{h,j}\Bigr)d_t+
 \sum_{\ell=t}^{t+k-1}
 \Bigl(\prod_{j=\ell+1}^{t+k-1}\rho_{h,j}\Bigr)b_\ell.
 \label{eq:cal-future-tube}
\end{equation}
对任务误差 $e=[I_m,-I_m]\chi$，相应管半径至多为 $\sqrt2d_j$；对纠偏速度，半径至多为 $\omega d_j$。能力验收可以使用整个预测状态域的 $R_{s,j}=\|\bar\chi_j\|_2+d_j$，或使用更紧的 $r_{\zeta,j}=\omega(\|\bar v_j\|_2+d_j)$。这给出沿名义计划滚动展开的未来管，而非仅围绕零误差的一条标量界。

近处有当前观测支持，远处累积响应变化、人体运输和目标漂移，因此经常呈现近窄远宽的形状；收缩强或噪声小时也可变窄，不能要求管宽永远单调增加。新的图像与探查更新初始集合、$X,\widehat J,q_j$ 及剩余任务参考，应从当前时刻重新递推\emph{整条}未来管。若要宣称新管是旧预测管条件化后的有效收紧，必须使用有效集合相交及一致模型；单纯替换 GP 均值不具有这种集合保证。

材料与姿态误差管经当前有效的身体映射运输到 CoP、姿态和机器人候选。局部映射若具有经过验证的 Lipschitz 常数 $L_P$，则材料位置半径 $d$ 给出物理位置半径至多 $L_Pd$，另加注册误差。GP 样本可生成候选未来形状与统计置信带，但未经覆盖率校准及模型验证的 GP 方差不能作为上述认证物理管半径。

\subsection{沿途有效区间与失锁模式}

一处探查不能保证整段长血管。若在共同运输标架中独立验证了
\begin{equation}
 \|\Theta(\xi)-\Theta(\xi_0)\|_2
      \le Ld_g(\xi,\xi_0),
 \label{eq:cal-spatial-variation}
\end{equation}
则当前 $N(\xi),C(\xi)$ 的证书加入
\begin{equation}
 q_{\mathrm{variation}}=
    Ld_g(\xi,\xi_0)\|N(\xi)C(\xi)\|_2.
 \label{eq:cal-variation-q}
\end{equation}
由此沿未来路径增大 $q_j$ 并重算整条管。新探查更新锚点可以收紧该项；调和坐标自身并不产生此变化界。没有有效变化界时，仅在局部窗口内提供保证。

目标丢失、耦合不足或关联不可信时，关闭图像误差驱动，推荐
\begin{equation}
 \dot\zeta=-2\omega\zeta,\qquad\dot s=0,
 \label{eq:cal-loss-mode}
\end{equation}
其中整个名义材料技能生成器随 $s$ 暂停，独立接触控制仍工作。这使纠偏速度逐渐制动，而不是保持最后一帧的非零纠偏速度。缺失检测不能填写成零误差。无观测时不使用定理~\ref{thm:cal-recovery}宣称任务恢复；图像特征出了视野甚至可能不再定义。

若存在单独有效的潜在目标位置预测模型，且固定实现的潜在任务响应满足 $\|JC\|_2\le M_C$、前馈及目标漂移界为 $w_{\mathrm{loss}}$，那么失锁后 $\tau$ 时间内的潜在误差半径可以保守扩张为
\begin{equation}
 d_e(t+\tau)\le d_e(t)
 +\frac{M_C\|\zeta(t)\|_2}{2\omega}
       (1-e^{-2\omega\tau})
 +w_{\mathrm{loss}}\tau.
 \label{eq:cal-loss-expansion}
\end{equation}
这是有条件的开环运动界，不是失锁图像的收敛证明。若没有这种模型，或预测超出可信窗口，就不能继续宣称拥有认证未来管。重获观测后首先验收目标关联及窗口，再运输状态、更新参考与重新修复。血管分叉需要保留互斥关联假设或报告歧义；单一高斯均值不能保证落在真实分支上。本文的响应校准定理不保证永不失锁，也不保证关联识别永远正确。

\section{有限候选算法的可实现与完备性边界}
\label{sec:cal-finite}

对每个保持材料任务的允许实现，生成有限且经过物理检查的探查库 $\{u_1,\ldots,u_M\}$。定义
\begin{equation}
  \mathcal L_N=
  \operatorname{span}\{Nu_1,\ldots,Nu_M\},
  \qquad r_{\mathrm{lib}}=\dim\mathcal L_N .
\end{equation}
实际算法依次执行探查并解~\eqref{eq:cal-convex-repair}。当验收失败时，可用当前候选的 $(I-P_X)NC$ 的主要左奇异方向提出下一条探查，再由实际能力检查和有限库选择决定是否执行。此贪心选择是启发式，不能声称最少探查最优。

在 $E_{\rm geo}=0$、无噪声且冻结的局部模型内，选择该库的一个实际可执行基，只需至多 $r_{\mathrm{lib}}$ 个独立方向。若某个实际可行任务实现 $C$ 满足
\begin{equation}
  JC=I_m,\qquad
  \operatorname{range}(NC)\subseteq\mathcal L_N ,
\end{equation}
则全部库基校准以后 $(I-P_X)NC=0$，其响应可由数据精确得到；凸可行性验收能够包含该实现。这个结论是对指定有限库和局部模型的条件完备性。噪声存在时，额外重复采样可能必要；它不能被独立方向秩的上界代替。

若没有独立安全探查可用，或允许实现均无法满足任务恢复证书，算法应明确返回不可验收，而不是继续执行一个只有名义可达性的方案。

\section{合法混合探查与后续实现可以相互适配}
以下是代数机制例，尚不是人体实验。取
\begin{equation}
 H=\begin{bmatrix}1&0&0\\0&1&0\end{bmatrix},\qquad
 N=\begin{bmatrix}1&0&1\\0&1&0\end{bmatrix}.
\end{equation}
任务零空间通道为 $\mathcal C=\operatorname{span}(1,0)$，最小曝光实现及其矩阵为
\begin{equation}
 B_\star=[e_1-e_3,e_2],\qquad
 D_\star=NB_\star=\begin{bmatrix}0&0\\0&1\end{bmatrix}.
\end{equation}
假设最便宜或唯一合法的探查为 $u_p=e_2+e_3$，则实际激励是 $q=Nu_p=(1,1)^\top$。若固定 $B_\star$，投影后曝光仍有
\begin{equation}
 \norm{(I-P_Q)D_\star}=1/\sqrt2,
 \qquad Q=\operatorname{span}(q).
\end{equation}
但 $\mathcal C+Q=\R^2$，正确剩余 $D_Q=0$。选择
\begin{equation}
 B_Q=[e_1-e_3,e_2+e_3],\quad
 HB_Q=I_2,\quad NB_Q=\begin{bmatrix}0&1\\0&1\end{bmatrix}
\end{equation}
即可把未来曝光放到已校准的混合方向。测到的向量 $g=\Theta q$ 给出
\begin{equation}
 JB_Q=I_2+[0,g],\qquad
 C_{\rm repair}=B_Q(I_2+[0,g])^{-1}.
\end{equation}
当该增益可逆且实际能力验收通过时，一次合法混合方向探查足够修复全部任务响应；若 $\delta\sqrt2<1$，可逆性对整个范数球成立。代价是 $\norm{B_Q}=\sqrt3$ 而 $\norm{B_\star}=\sqrt2$。

这里不要求 IRD 产生恰好等于理想 SVD 的动作。实际探查可以受约束；机器人也可以改变后续实现，使未来所需反馈落入真实已经校准的方向。这一双向匹配由 $D_Q$ 定理描述，而不是额外增加一个学习模块。

\section{超声中的真实不确定通道：由切片交点直接推导}
\subsection{面内动作和面外动作需要不同的未知量}
在有效、关联正确的局部窗口，血管切向为 $t$，探头原点为 $p$，切片法向为 $n$，面内单位基为 $E=[e_x,e_z]$。局部交点由
\begin{equation}
 \gamma(s)=\gamma_0+ts,\qquad
 n^\top(\gamma(s)-p)=0,\qquad
 y=E^\top(\gamma(s)-p)
\end{equation}
给出。固定姿态对时间微分，若 $n^\top t\ne0$，则
\begin{align}
 \dot s&=\frac{n^\top v}{n^\top t},\\
 \dot y&=-E^\top v+\beta n^\top v,
 &\beta&=\frac{E^\top t}{n^\top t}.
\end{align}
二维单切片看不到完整面外血管切向，因此 $\beta$ 需要先验或校准。若 $n^\top t$ 太小，交点模型病态，应拒绝该窗口而非宣布小样本可修复。绕切片法向的转角符号取 $\dot y=(z,-x)^\top\dot\alpha$，该动作不改变切片本身。

局部平面皮肤下，选择面内切向平移、绕切片法向摇摆、面外皮肤切向平移三种控制，得到
\begin{equation}
 J(y)=
 \begin{bmatrix}-1&z&\beta_x\\0&-x&\beta_z\end{bmatrix}
 =H+\theta N,\quad
 H=\begin{bmatrix}-1&z&\bar\beta_x\\0&-x&\bar\beta_z\end{bmatrix},\quad
 N=\begin{bmatrix}0&0&1\end{bmatrix}.
\end{equation}
这里 $\theta=\beta-\bar\beta\in\R^{2\times1}$，不是人为添加的低秩未知模型：面外斜率只在面外位移时进入，具有实际几何来源。若两分量均未知，任意有限 $\beta$ 都对应真实单位切向
\begin{equation}
 t(\beta)=\frac{E\beta+n}{\sqrt{1+\norm{\beta}^2}},
\end{equation}
因此局部斜率球可作为明示的物理模型类。若另有解剖约束使未知量只有一个标量或更小集合，谱范数球只是保守包络：修复上界仍有效，精确最小最大下界不能直接宣称在缩小后的物理类中仍紧。未知的接触变形、传感标定误差、关联错误不能由这个因式分解自动消除；前三者的可界定部分进入剩余项，关联错误则使模型无效。

中心实现 $x=0$、$\bar\beta_z\ne0$ 时，$H$ 满行秩，但 $H\ker N$ 只能产生一维面内任务作用。因此局部首次单位响应的结构校准方向数为一。若任务允许 $x=\ell\ne0$，则
\begin{equation}
 B_\ell=
 \begin{bmatrix}
 -1&-z/\ell\\0&-1/\ell\\0&0
 \end{bmatrix},\qquad
 JB_\ell=I_2,\qquad NB_\ell=0.
\end{equation}
两维任务仍被控制，而未知 $\beta$ 不被调用。这不是被动换坐标，而是换了实际机器人实现与真实杠杆。另一方面，$1/|\ell|$ 放大了运动需求；太小的偏移可能违反角速度、接触和关节能力约束。

\subsection{有限动作必须用精确接触几何验收}
固定切片平面、静止目标交点和局部平面皮肤；CoP 沿固定皮肤切向平移 $d$，探头在该平面中转动 $\alpha$。则
\begin{equation}
 y'=R_2(-\alpha)\begin{bmatrix}x-d\\z\end{bmatrix},\qquad
 R_2(\alpha)=\begin{bmatrix}\cos\alpha&-\sin\alpha\\\sin\alpha&\cos\alpha\end{bmatrix}.
\end{equation}
目标 $y'=(\ell,z_\star)$ 等价于
\begin{align}
 z&=\ell\sin\alpha+z_\star\cos\alpha,\\
 d&=x-\ell\cos\alpha+z_\star\sin\alpha.
\end{align}
第一式在零角附近的导数为 $\ell$；$\ell\ne0$ 时由隐函数定理得到局部小角解。无角度限制下仍有正深度偏差硬界
\begin{equation}
 z-z_\star\le\sqrt{\ell^2+z_\star^2}-z_\star.
\end{equation}
若有 $|\alpha|\le\alpha_{\max}$，需要取该区间上 $\ell\sin\alpha+z_\star\cos\alpha$ 的实际极值。不能只用一阶 Jacobian 宣称任意有限动作可达。

全路径接触约束下的已知交互矩阵为
\begin{equation}
 H_0(\alpha)=\begin{bmatrix}-\cos\alpha&z'\\\sin\alpha&-x'\end{bmatrix},\qquad
 \det H_0=x'\cos\alpha-z'\sin\alpha=x-d.
\end{equation}
持续可控需要物理横向杠杆 $|x-d|$ 全程有下界；图像横向偏移本身不是完整条件。有限控制可以直接求解上述精确几何，避免把已知的大转角非线性误差塞成不必要的校准。

\noindent\textbf{纯几何核验。}取 $x=\ell=5\,\mathrm{mm}$、$z_\star=30\,\mathrm{mm}$，正深度硬界为 $0.41381265\,\mathrm{mm}$。$\pm0.5\,\mathrm{mm}$ 的对称修复不成立。对 $\pm0.1\,\mathrm{mm}$，精确解为
\begin{center}
\begin{tabular}{rrrr}
\toprule
初始深度差 & 转角 & CoP 平移 & 终点目标误差\\
\midrule
$-0.1\,\mathrm{mm}$ & $-1.08441^\circ$ & $-0.566866\,\mathrm{mm}$ & $<10^{-12}\,\mathrm{mm}$\\
$+0.1\,\mathrm{mm}$ & $+1.22452^\circ$ & $+0.642248\,\mathrm{mm}$ & $<10^{-12}\,\mathrm{mm}$\\
\bottomrule
\end{tabular}
\end{center}
这些是精确刚体模型的数值核验，不是真实探头实验。声学耦合、压力、组织变形、IK 和碰撞仍需独立验收。

\subsection{同一任务、首次合同与总成本的边界}
面外移动会改变交点的材料相位，故共同任务必须允许同一血管片段内的有限相位偏差。若要求严格固定材料截面，则不能将面外中心纠偏与固定平面偏心纠偏当成同一任务。

中心时的一方向下界只对应冻结局部模型与首次合同。任意多步策略可以先平移形成杠杆再摇摆，因此不能声称所有从中心出发的扫描策略都必须校准 $\beta$。这条重构路径本身也是可选实现，其耗时必须计入。

长血管的推进通常仍含面外前馈 $u_{\rm ff}$，产生未知漂移 $\Theta Nu_{\rm ff}$；它进入恢复定理中的 $\widetilde w$，例如需要纳入 $\delta\norm{Nu_{\rm ff}}$ 的界。$NB_\ell=0$ 免去的是这两个局部特征纠偏的斜率校准，不是免去全扫描的所有未知量。若任务合同还要求精确控制推进相位或额外切面特征，应把这些要求加入任务模型，重新计算 $H,N,D_\star$，不能搬用二维结论。

按横向速度上限 $15\,\mathrm{mm/s}$，从中心转到 $5\,\mathrm{mm}$ 偏移至少要 $0.333\,\mathrm{s}$。免去一个校准方向不保证总时间降低；只有
\begin{equation}
 T_{\rm switch}+T_{\rm probe}+T_{\rm repair}+T_{\rm execute}
\end{equation}
在同一任务质量和真实预算下更小，才构成工程收益。

\section{可微 IRD 怎样进入完整算法，而不膨胀系统}
\subsection{候选生成与硬验收分开}
令 $a$ 包括材料允许偏离、探头姿态纤维、rail 配置和机器人实现分支。变换到 IRD 的完整探头目标为
\begin{equation}
 T_{\rm rel}(a)=T_{\rm base}(a)^{-1}T_{\rm probe}(a).
\end{equation}
对可微 IRD 得分 $\mathcal I(T_{\rm rel},a)$，用链式法则
\begin{equation}
 D_a\mathcal I=D_T\mathcal I\,D_aT_{\rm rel}+\partial_a\mathcal I
\end{equation}
搜索能保留任务且具有足够纠偏能力的候选。校准尾谱可作目标项：
\begin{equation}
 \min_{a\in\mathcal A_\tau}
 c_{\rm switch}(a)+c_{\rm probe}(a,k)
 +\lambda\,\sigma_{k+1}(D_\star(a))
 -\mu\,\mathcal I(T_{\rm rel}(a),a).
\end{equation}
奇异值重合处用次梯度、光滑近似或直接枚举小候选库。IRD 是提议器；接受条件来自完整 IK 分支、接触、多面体输入能力、碰撞和有效观测，而非一个软得分。选姿态时 $H,N$ 也可能改变，必须逐候选重算，不搬用旧的 $D_\star$。

\subsection{有限候选库中的可审查算法}
\begin{enumerate}
\item 由材料任务及纤维记录运输任务目的，建立当前允许实现集合；检查动捕、图册与可信目标关联。
\item 由实际几何计算 $H_a,N_a$ 与剩余误差界，先纳入接触约束，再白化度量。
\item 对每个真实可行候选计算 $D_\star(a)$，由谱尾得到必要有效激励方向数的放宽下界。
\item 生成可执行的短探查库。每条候选存储实际可激励通道、时间、幅度、接触变化和观测有效区间；不是只存一列理想 SVD 向量。
\item 在库内枚举或分支定界探查组合；使用预期噪声界与激励条件排除无法达成目标性能的组合。这里的最优性仅相对于声明的有限库。
\item 执行选定短探查，记录实际材料拉回动作。失锁段不能作为可信特征回归数据；模型有效性被破坏时重新获取目标。
\item 用实际数据求解修复实现 $C$ 的凸验收问题。证书不通过时，补探查、换实现或停止推进；不能把不可行结果剪掉后声称完备。
\item 执行闭环短段，按沿途变化界更新未来管，并在证书到期前重算或探查。
\end{enumerate}
对每个候选得到可行总代价上界 $U_a$。无约束谱给出方向下界；若每有效方向确有统一代价下界 $c_{\min}>0$，可转为 $L_a=c_{\rm switch}^{\rm lb}+c_{\min}k_{\rm necessary}$。没有实际时间或能量下界时，不能把方向数直接当作时间下界。选出的 $U$ 与放宽 $L$ 构成可报告差距；它不是整个连续非线性扫描问题的全球最优证明。

\section{GP 几何技能、IRD 与最小反馈校准的完整组合}
\label{sec:v43:integration}
\subsection{核心保持不变，新增生成器只提出技能偏好}
GP 输出在材料技能空间中有物理语义的 forcing，进度一致的几何二阶系统生成 $Y_{\rm pref}$；投影规则将内部任务、表面作用点与姿态纤维联系起来。对当前数据 $X$ 与候选实现 $a$，联合选择采用
\begin{equation}
\begin{aligned}
 \min_{a,C,\nu}\quad &c_{\rm switch}(a)+c_{\rm probe}(a,X)
       +\lambda_{\rm style} J_{\rm style}(a)-\lambda_I\mathcal I(T_{\rm rel}(a),a)
       +\lambda_u\|C\|_F^2,\\
 \text{s.t.}\quad &a\in\mathcal A_\tau,\qquad \widehat J_aC=I_m,\qquad q_a(C,X)\le q_{\rm goal}<1,\\
 &E_aC=0,\qquad u_{\rm ff}(a,\nu)+C\zeta\in\mathcal U_a\quad\forall\|\zeta\|\le r_\zeta,\\
 &\text{有效目标关联、接触、完整 IK 分支、碰撞及速度变化界通过验收}.
\end{aligned}
\label{eq:v43:joint-choice}
\end{equation}
联合选择并非凸问题。仅在固定 $a$、数据、前馈和可信速度半径时，响应修复子问题是凸的；dIRD 梯度负责候选提议，有限候选枚举与硬验收给出审查边界。谱尾用于说明还缺多少独立校准方向，不把一个低 IRD 代价当作患者反馈已经校准。

\begin{proposition}[技能生成表达不改变局部校准核心]
\label{prop:v43:core-unchanged}
固定相同材料任务维数、允许物理实现、度量及局部响应类。将固定轨迹、GP 参考或 GP-forcing 几何生成器用作名义偏好，不改变 $D_\star,D_Q$ 的定义与其首次任务下界。对一个被接受的实现，只要由实际数据获得 $\widehat JC=I$、$q<1$，且生成的前馈残差及控制执行误差满足已声明扰动界，二阶恢复合同仍成立。
\end{proposition}
\begin{proof}
首次下界仅由局部 $H,N$、任务零空间及实际信息子空间 $Q$ 定义；名义偏好不是这些对象的额外观测。生成器可以影响候选 $a$，故每个被选择候选须重算 $H_a,N_a,D_Q(a)$。接受以后，将 $Ju_{\rm ff}-\dot y_r$ 与执行残差放入 $\widetilde w$，即可使用二阶恢复定理。该推导没有将名义生成器稳定性转化为未经校准的真实反馈稳定性。
\end{proof}

这给出系统连接的具体含义：调和图册与纤维保留可迁移任务及相对技能；GP 学偏好；几何二阶生成器实现形状和状态接续；IRD 利用允许偏离改变物理实现；商空间谱说明该实现仍必须知道哪些响应；真实探查与凸修复最后给出反馈性能验收。它们的职责相接，但任何一个模块都不能替代另一个模块的证据。

\subsection{校准质量直接约束材料技能的允许推进速度}
这一接口使 GP forcing、IRD 和校准证书在一个可计算量上相接。由几何技能与实现选择得到每单位材料进度所需的真实动作 $b_s$，令 $u_{\rm ff}=\nu b_s$，并在扣除已知人体运动的任务模型内满足 $Hb_s=y_r'(s)$。将其它保持运动与外部漂移放入 $w_{\rm ext}$，则
\begin{equation}
 \widetilde w=\nu\Theta Nb_s+\nu E_{\rm geo}b_s+w_{\rm ext},\qquad
 \|\widetilde w\|\le\nu\alpha_s+w_{{\rm ext},\max},\quad
 \alpha_s=\delta\|Nb_s\|+\eta\|b_s\|.
\end{equation}
这里 $Hb_s=y_r'$ 是前馈兼容性的验收条件，不由 GP forcing 自动保证。若存在已知名义不匹配，$\alpha_s$ 应再加 $\|Hb_s-y_r'\|$。若用 $\widehat Jb_s=y_r'$ 构造前馈，可以用实际探查数据给出的方向误差界替换该保守 $\alpha_s$，但必须重新验收。令 $\lambda_s=\omega(1-q)/2$，二阶稳态残差管的充分任务条件为
\begin{equation}
 \frac{\sqrt2\,[\nu\alpha_s+w_{{\rm ext},\max}+\sqrt2\omega n_{\max}]}{\lambda_s}
 \le\epsilon_{\rm task}.
\label{eq:v43:progress-contract}
\end{equation}
若还要求整个执行段 $\|e(t)\|\le\epsilon_{\rm task}$，必须同时有 $S(t_0)\le\epsilon_{\rm task}/\sqrt2$，或逐时验收完整瞬态管。仅降低推进速度不能让已有大误差瞬间消失。当 $\alpha_s>0$ 时，上式给出稳态合同下允许的推进上限
\begin{equation}
 \nu\le\nu_{\rm cert}=
 \frac{\lambda_s\epsilon_{\rm task}/\sqrt2-w_{{\rm ext},\max}-\sqrt2\omega n_{\max}}{\alpha_s}.
\end{equation}
分子为负时不存在非负推进速度满足该充分合同，应先改善观测、校准或实现，而不是把负值裁成零后宣布已通过。$\alpha_s=0$ 时由剩余条件单独判定。实际 $\nu$ 还需同时满足 IRD 候选经硬验收的输入能力、参考物理速度和 $\dot\nu Z+\nu^2\nabla_s Z$ 加速度界。该上限为充分而非最优速度；它说明技能越复杂或前馈曝光越大，所需校准与剩余能力越重要。即使局部纠偏满足 $NC=0$，沿血管推进的 $Nb_s$ 也可能非零，不能因此免去全部扫描不确定性。

\subsection{最小运行状态与闭环图}
在一个可信短窗口，实际保存的状态为：部位与目标分支标签、材料进度 $s$、名义技能 $(Y,Z)$、二阶图像纠偏状态 $\zeta$、实现 $a$、实际探查数据 $(X,R_{\rm data})$、修复 $C$、响应可信界及未来管。GP 超参数和后验在低速更新层维护；不新增 KMP、另一套运动基元学习器或图像基础模型。

\begin{figure}[htbp]
\centering
\begin{tikzpicture}[node distance=9mm and 10mm,>=Latex,
 box/.style={draw=blue!45!black,rounded corners,fill=blue!4,align=center,text width=56mm,minimum height=12mm,font=\small},
 arrow/.style={->,thick,draw=blue!50!black}]
 \node[box] (skill) {材料任务、内部先验\\GP forcing 与几何技能状态};
 \node[box,right=of skill] (ird) {允许实现与 dIRD 提议\\接触、IK、碰撞验收};
 \node[box,below=of skill] (data) {实际动作与有效图像\\身体运动补偿、探查数据};
 \node[box,below=of ird] (repair) {任务商空间谱与凸修复\\$D_Q$, $\widehat J C=I$, $q<1$};
 \node[box,below=of data] (servo) {名义前馈与二阶纠偏\\$u=u_{\rm ff}+C\zeta$};
 \node[box,below=of repair] (tube) {增广未来管与模式判断\\推进、补探查、制动重捕获};
 \draw[arrow] (skill)--(ird);
 \draw[arrow] (ird)--(repair);
 \draw[arrow] (data)--(repair);
 \draw[arrow] (repair)--(tube);
 \draw[arrow] (repair)--(servo);
 \draw[arrow] (servo)--(data);
 \draw[arrow] (tube)--(servo);
 \draw[arrow] (tube.east)--++(6mm,0)|-(ird.east);
\end{tikzpicture}
\caption{只增加明确的技能生成状态，不改变校准决策核心。图像失锁时停止证据更新，名义进度暂停，接触环与纠偏制动仍继续。}
\end{figure}

\subsection{数据复用必须区分被动换图与主动改模型}
实际探查数据 $X_a,R_{{\rm data},a}$ 属于所声明实现与局部响应模型。被动换图时可按已知算子精确运输；主动改变切片法向、接触状态或物理姿态后，$\Theta_a$ 不一定等于旧 $\Theta$ 的已知运输。只有通道对应与误差变化界通过验证时才复用旧数据，否则维护候选模型索引或重新探查，并将有界变化纳入 $\eta$、$q_{\rm variation}$。不能让所有候选无条件共享同一个 GP 后验和校准矩阵。改变纯控制分配 $B$ 而不改变同一个 $J,H,N$，才是前文 $D_Q$ 中可直接复用同一信息子空间的情形。

\subsection{物理速度重建与当前状态接续}
归一化材料进度不等于相同物理扫描速度。对新身体中接受的非退化作用点曲线 $\xi(s)$，若示教物理速度偏好为 $v_{\rm pref}(s)$，则
\begin{equation}
 \nu_{\rm pref}=\frac{v_{\rm pref}(s)}{\|D_\xi P\,\xi'(s)\|},\qquad
 v_{\rm CoP}=D_\xi P\,\xi'(s)\nu+b_p.
\end{equation}
实际 $\nu$ 再受式~\eqref{eq:v43:progress-contract}及机器人能力限制；曲线导数接近零时转入停留或姿态模式，不作该除法。压力、姿态和速度偏好是不同的示教量，不能由速度归一化推导出接触阻抗。

实际物理加速度还包括身体映射的曲率及时间变化项：
\begin{equation}
 a_{\rm CoP}=J_P\ddot\xi+\dot J_P\dot\xi+\dot b_p.
\end{equation}
$\dot J_P$ 含沿材料运动的 Hessian 作用和人体、接触状态的变化；$\dot b_p$ 也需要相应估计或可信界。不能仅验收材料生成器中的 $\dot\nu z_b+\nu^2D_sz_b$ 就宣称物理加速度已保证。

切换图、分支或几何参考时，在局部可逆表示 $Y=\Phi_{\rm new}(t,s,x)$ 中，可以按旧参考状态 $Y_o,\dot Y_o$ 初始化
\begin{equation}
 x_n=\Phi_{\rm new}^{-1}(Y_o),\qquad
 D_x\Phi_{\rm new}\,\dot x_n
 =\dot Y_o-\partial_s\Phi_{\rm new}\,\nu-\partial_t\Phi_{\rm new}.
\end{equation}
等式中的向量先运输到同一切空间。姿态参考 $R_o$ 的具体匹配为
\begin{equation}
 \Upsilon_n=R_{0,n}^\top R_o,\qquad
 \omega_{\Upsilon,n}=\omega_{R_o}-\Upsilon_n^\top\omega_{0,n}.
\end{equation}
在 $\nu>0$ 时，时间速度还可转换为空间生成速度；在暂停时不除以 $\nu$，先保持当前材料状态，按有限短段桥接后恢复推进。这个初始化只在映射有效、速度匹配可行且接续后的能力及初始增广状态域验收通过时提供连续性。无法匹配时采用显式有限桥接或停止，不宣称任意目标切换自动达到 $C^1$ 或 $C^2$。实际 $C,u_{\rm ff}$ 的接续也要避免命令突变，并单独验收。

\subsection{可执行算法：提议、探查、修复与接续}
\begin{enumerate}
 \item 将当前材料任务、目标分支与 GP 技能后验运输到有效图。检查周期角 lift、旋转对数域及动捕可信区间。
 \item 在材料进度上积分名义技能；从内部几何与残余姿态构造作用点和探头姿态，计算物理速度及 $\dot\nu Z$ 加速度项。该结果是偏好，不锁定唯一世界 TCP 轨迹。
 \item dIRD 结合技能代价与校准谱提出少量允许实现；逐个重算 $H_a,N_a$，纳入接触约束并在固定物理度量下白化。
 \item 使用已有实际数据试解修复。若可行，检查二阶状态包络下的输入能力、前馈残差与短窗口观测条件。
 \item 若不可行，在有限安全探查库中选择能补足 $(I-P_X)N_aC$ 主要方向的短动作；切换实现与探查的时间和扰动一并计入预算。
 \item 仅从关联正确、质量合格的观测和实际材料动作更新 $X,R_{\rm data}$。名义命令与实际执行不同的部分要记录，而不能把名义命令当真值。
 \item 求解响应修复，按增广二阶证书重算整段未来管；通过才推进。失败则补探查、换允许实现或停止。
 \item 失锁时平滑令 $\nu\to0$，完整暂停名义技能，并制动 $\zeta$；接触环继续。重捕获后先验证目标关联和旧证书适用性，再按当前状态接续。
\end{enumerate}
候选库穷尽后的失败是“在声明的库和模型内不可验收”，不是原任务必然不可完成。计算预算耗尽、秩变化、图册失效或模型界失效也应作为明确返回状态。

\subsection{未来位置锥不要求概率，但需要正确的证据语义}
由二阶增广状态管得到任务误差集合；若任务到位姿的局部映射 $\Phi_a$ 在可信域满足 Lipschitz 界 $L_\Phi$，则由误差引起的物理位置偏差可包络为 $L_\Phi\sqrt2 R$。完整未来集合还要加上内部解剖位置集合、名义技能候选变化及人体运输误差：
\begin{equation}
 \mathcal P_{t:j}=\bigcup_{b,a}\Phi_{a,b}\bigl(\mathcal A_{{\rm anat},b,j},\mathcal Y_{{\rm skill},j}\bigr)
 \oplus\mathcal B\bigl(L_\Phi\sqrt2 R_j+\epsilon_{{\rm map},j}\bigr).
\end{equation}
这里的 $\Phi_a$ 必须是经过验收的实际局部实现映射；$m<6$ 的超声特征误差并不约束全部位姿零空间，$HB=I$ 也不自动给出可积逆映射。不存在该映射时只认证任务特征管，物理未来管须由完整位姿状态的前向传播另行得到，不能直接使用 $L_\Phi\sqrt2 R$。
并集保留分叉；近未来可能因为刚校准和观测有效而较窄，远未来因响应变化界、关联和几何误差而变宽。新数据会更新锚点、候选和所有 $R_j$。若集合跨过图边界或 Log 域，应拆成多图候选，不能用一个无限扩大的局部高斯继续解释。

GP 后验协方差可生成条件于统计模型的解剖和技能候选区间；它不是 $\delta,\eta,L$ 的确定性证书。给概率覆盖率时必须报告模型、校准方法及验证数据。此处的鲁棒恢复管则条件于明示界，两类不确定性可以在预测中相加，但不能混写成一个未经验证的概率安全保证。

\section{贡献定位、验证设计与当前证据}
\label{sec:v43:assessment}
\subsection{唯一主贡献及已有基础的边界}
本版主张保持为：在材料任务、观测质量及允许偏离事先固定时，通过选择真实机器人实现降低不可避免的反馈曝光；用实际可执行探查补足必要方向，并将校准质量转为任务恢复、输入能力与允许推进速度的共同合同。最强理论对象仍是 $D_Q=P_{\mathcal C+Q}^{\perp}NH^\dagger$，而不是新增一种运动基元。

任务导向探索已有任务相关样本复杂度\cite{wagenmaker2021}，最低输入设计已有任务属性辨识\cite{kang2023}；不确定控制分配已考虑效能模型的参数不确定性\cite{grauer2022}，性能导向探查已有顺序学习与控制\cite{venkata2024}，部分信息与最坏恢复也有经典基础\cite{foucart2024}。任务兼容主动校准还强调探查后的终端可执行性\cite{kobayashi2026}。本稿的区别限于明确模型类内的共同商空间、首次任意物理动作的精确下界、谱预算前沿，以及合法混合探查和后续实现的双向适配；不声称首次发现任务相关辨识。

GP 学 DMP forcing\cite{gpdmp2026}、弧长 DMP\cite{aldmp2016}、几何 DMP\cite{gdmp2024}和流形向量 GP\cite{gpmanifold2021}均是已有基础。4.3 的新增价值是把这些基础接到真实材料坐标、图册转换、允许实现和同一响应证书上，并修正二阶能力与未来管。它们提高系统完整性，不自动提高已有部件本身的原创性。超声交互矩阵与自适应视觉伺服同样已有先例\cite{krupa2005}。

\subsection{最强对照必须拥有相同自由度}
主对照是相同目标提取、观测模型、接触模型、允许偏离、实现分支及实际探查库下的 belief-space/dual MPC。该对照必须也能改变偏心目标、姿态、rail 与探查动作，并使用相同时间和计算预算；本方法不独占有利动作。还应配置充足计算的离线参照，以区分决策近似与执行效果，但不能把离线预算用于在线公平主对照。

在相同材料坐标、数据和执行带宽下，对比 GP 参考加二阶跟踪与 GP forcing 加几何空间生成。固定调度下二者可能通过线性微分算子等价，因此检验重点是变速、暂停接续、姿态图转换和少量示教下的形状保持。笛卡尔 GP 应有相同人体跟踪补偿和在线反馈，不能故意使其失去本来可用的信息。

主要消融为：只优化 IRD 能力；固定实现后校准；有响应不确定分配但无性能驱动探查；有探查但不允许任务等价实现；去掉 forcing 的几何运输；去掉增广二阶证书。报告同一合同下的有效谱尾、真实独立激励方向、实际探查与切换时间、恢复率、误差管覆盖、有效声窗、接触变化、总代价和总计算时间。方向数下降但总时间增加，或 $E_{\rm geo}$ 支配收益时，应报告负结果。

\subsection{合成数据与少量真人数据分别能支持什么}
合成 SMPL-X 人体、内部血管分支和器官区域可以产生明确材料真值，用于训练或检查解剖趋势、几何投影、技能运输以及候选生成。它不能提供真实超声分割误差、组织接触形变、医生动作残余或压力手感。专家动捕与同步探头观测提供技能 forcing；实际短探查提供患者局部响应信息；独立留出人体用于验证模型界和泛化。

可以先利用少量示教限制在低维、短片段、已定义任务上下文的 GP 技能类；无需以大规模人类数据作为理论前提。但样本数应由留出人体的误差与覆盖率确定，不能事先保证若干条示教就足够。若没有测力数据，只报告运动与姿态形状的保持，不报告已学习完整触觉手感。每个患者仍需要实际观测与证书验收，不能由合成数据直接替代。

\subsection{可复现数学与几何核验}
本次使用固定随机种子 $20261009$，执行随源文件附带的 \texttt{validate\_v43.py}。表中结果只检查数学实现及几何一致性，不是机器人或真人扫描实验。
\begin{center}
\small
\begin{tabularx}{\linewidth}{Yl}
\toprule
核验对象 & 实际结果\\
\midrule
75 组矩阵、750 个右逆：$HB_\star=I$ & 最大误差 $1.59\times10^{-14}$\\
最小曝光 Gram 恒等式 & 最大误差 $8.91\times10^{-14}$\\
合法混合探查 $D_Q$ / 修复响应 & $2.51\times10^{-16}$ / $1.14\times10^{-16}$\\
非恒定正交标架下 GP 物理均值 / 协方差 & $2.14\times10^{-14}$ / $1.29\times10^{-15}$\\
协变 forcing 标签 & 最大误差 $5.97\times10^{-15}$\\
遗漏 connection 的负对照 & 标签误差 $5.106$\\
AL 双状态含 $2\,\mathrm{s}$ 暂停 & 形状误差 $6.60\times10^{-11}$\\
二阶恢复 Lyapunov 随机核验 & $10{,}000$ 例无证书违例\\
超声有限偏心刚体几何 & 终点误差 $\le3.67\times10^{-15}\,\mathrm{mm}$\\
\bottomrule
\end{tabularx}
\end{center}
GP 核、均值、标签和各向异性观测噪声全部按同一标架运输。遗漏噪声运输的负对照产生 $0.01299$ 的物理均值误差。AL 核验使用解析 forcing，GP 后验核验使用独立合成标签；没有把这两个核验拼成未经执行的 GP 学习机器人实验。随机核验支持代码一致性，不能替代前文证明。若使用随包提供的图，其中 GP 的 $\pm2$ 标准差带明确未做覆盖率校准。

\IfFileExists{v43_validation.pdf}{
\begin{figure}[htbp]
\centering\includegraphics[width=\linewidth]{v43_validation.pdf}
\caption{可复现核验：GP 标架一致性、AL 暂停与变速、二阶恢复包络。均为声明模型中的合成数学测试，非真实任务性能比较。}
\end{figure}
}{}

\subsection{独立审查评价的范围}
两名独立审读代理分别检查数学核心与系统必要接口，并对完整 4.3 做第二轮复审。最终共同认可：在明确局部模型类内，商空间首次任务下界与实际探查--后续实现匹配的核心原创性、理论深度各为 $4/5$。数学审读对系统一致性的 $4/5$ 评价要求落实局部域、前馈兼容、初始包络、两阶段 GP 和位姿实现映射等边界；本稿已补入这些条件。系统审读要求的一般式余项、实际输入紧化及过宽公式亦已修订。

\begin{center}\small
\begin{tabularx}{\linewidth}{Ycc}
\toprule
评价对象 & 数学审读 & 系统审读\\
\midrule
核心原创性（局部模型类） & $4/5$ & $4/5$\\
理论深度 & $4/5$ & $4/5$\\
系统数学一致性／接口完整性 & $4/5$ & $4/5$\\
工程合理性（研究设计） & 未单独评分 & $4/5$\\
解释直观性 & 未单独评分 & $4/5$\\
真实任务验证成熟度 & 不足 $4/5$ & $2/5$\\
\bottomrule
\end{tabularx}\end{center}

这不表示所有单个模块都具有 $4/5$ 原创性，也不表示已有整篇论文的验证质量达到 $4/5$。当前工程合理性可以评为有依据的研究设计；真人与机器人结果尚未提供，验证成熟度仍不足。评分是独立代理的研究意见，不能作为外部同行评审、投稿保证或新增实验结果。完整结论应围绕理论收益与适用边界，而不是为满足评分而继续增加模块。

\subsection{最终取舍}
唯一技能学习路线为 GP forcing；学习在材料切空间与相对姿态纤维中，世界笛卡尔坐标用于重建和真实执行。空间二阶生成器负责几何形状与进度一致性，固定增益二阶图像纠偏负责有效观测下的任务恢复。IRD 仍选择允许实现，真实探查及商空间谱仍决定必须校准的方向，凸修复仍提供实际响应证书。这一选择保持核心 idea，并将各项接口补全；是否带来真实总成本和迁移收益，必须由同预算强对照验证。

\begin{thebibliography}{99}
\bibitem{pavlakos2019} G. Pavlakos et al. Expressive Body Capture: 3D Hands, Face, and Body from a Single Image. CVPR, 2019. \url{https://arxiv.org/abs/1904.05866}
\bibitem{wagenmaker2021} A. Wagenmaker, M. Simchowitz, K. Jamieson. Task-Optimal Exploration in Linear Dynamical Systems. ICML, PMLR 139:10641--10652, 2021. \url{https://proceedings.mlr.press/v139/wagenmaker21a.html}
\bibitem{kang2023} S. Kang, K. You. Minimum input design for direct data-driven property identification of unknown linear systems. Automatica 156:111130, 2023. \url{https://doi.org/10.1016/j.automatica.2023.111130}
\bibitem{grauer2022} J. A. Grauer, J. Pei. Minimum-Variance Control Allocation Considering Parametric Model Uncertainty. AIAA SciTech, 2022. NASA NTRS 20210023819. \url{https://ntrs.nasa.gov/citations/20210023819}
\bibitem{venkata2024} J. Venkatasubramanian, J. K\"ohler, J. Berberich, F. Allg\"ower. Sequential learning and control: Targeted exploration for robust performance. arXiv:2301.07995, version 3, 2024. \url{https://arxiv.org/abs/2301.07995}
\bibitem{kobayashi2026} J. Kobayashi. When Identifiability Is Not Enough: Terminal Controllability in Task-Compatible Active Calibration. IEEE Robotics and Automation Letters 11(10):11039--11045, 2026. \url{https://doi.org/10.1109/LRA.2026.3719194}
\bibitem{krupa2005} A. Krupa, F. Chaumette. Control of an ultrasound probe by adaptive visual servoing. IROS, 2005. \url{https://doi.org/10.1109/IROS.2005.1545272}
\bibitem{foucart2024} S. Foucart, C. Liao. Radius of Information for Two Intersected Centered Hyperellipsoids and Implications in Optimal Recovery from Inaccurate Data. arXiv:2401.11112, 2024. \url{https://arxiv.org/abs/2401.11112}
\bibitem{aldmp2016} A. Ude, R. Vuga, B. Nemec, J. Morimoto. Trajectory Representation by Nonlinear Scaling of Dynamic Movement Primitives. IROS, 2016. \url{https://hcr.ijs.si/resources/papers/iros16.pdf}
\bibitem{gdmp2024} F. J. Abu-Dakka, M. Saveriano, V. Kyrki. A unified formulation of geometry-aware discrete dynamic movement primitives. Neurocomputing 598:128056, 2024. \url{https://doi.org/10.1016/j.neucom.2024.128056}
\bibitem{gpdmp2026} W. Si, Z. Jin, N. Wang, C. Yang. Robot Learning System Based on Target Localization and Human Demonstration for Medical Examination. IEEE/ASME Transactions on Mechatronics, accepted manuscript, 2026. \url{https://repository.essex.ac.uk/42960/1/Robot_Learning_System_Based_on_Target_Localization_and_Human_Demonstration_for_Medical_Examination.pdf}
\bibitem{gpmanifold2021} M. Hutchinson, A. Terenin, V. Borovitskiy, S. Takao, Y. W. Teh, M. P. Deisenroth. Vector-valued Gaussian Processes on Riemannian Manifolds via Gauge Independent Projected Kernels. NeurIPS, 2021. \url{https://arxiv.org/abs/2110.14423}
\end{thebibliography}
\end{document}
